\documentclass[12pt,openany]{book}

% 1. Layout & Graphics
\usepackage[inner=20mm,outer=15mm,top=20mm,bottom=15mm]{geometry}
\usepackage{wrapfig}
\usepackage{caption}
\usepackage{graphicx}
\usepackage[export]{adjustbox}
%\graphicspath{{./images/}}

% 2. Language settings (for XeLaTeX and Pandoc)
\usepackage{iftex}
\ifxetex
\usepackage{fontspec}
\usepackage{polyglossia}
\setmainlanguage{malayalam}
\setmainfont[Script=Malayalam, HyphenChar="00AD]{RIT Rachana}
\setlanghyphenmins{malayalam}{3}{4}
\linespread{1.2}
\widowpenalty=10000
\clubpenalty=10000
\raggedbottom
\sloppy
\else
\usepackage[utf8]{inputenc}
\usepackage[T1]{fontenc}
\fi

% 3. Text formatting
\usepackage{anyfontsize}
\usepackage{setspace}
\usepackage{color}
\usepackage{quoting}
\usepackage{pdfpages}

% 4. Custom environment for verses
\newenvironment{myverse}
{\begin{quoting}[leftmargin=3em, rightmargin=3em, vskip=2ex, indentfirst=false, font=small]}
	{\end{quoting}}

% 5. Hyperlinks
\usepackage[colorlinks=true, urlcolor=blue, linkcolor=red]{hyperref}

% To reduce unnecessary gaps
\setlength{\abovecaptionskip}{-2pt} 

\pagestyle{plain}
\title{\textbf{ശ്രീ നാരായണ ഗുരു}}
\author{കുമാരനാശാൻ}
\date{}

\begin{document}
	
	\includepdf[pages={1}, viewport=0 0 595 842, scale=1.0, fitpaper=true]{guru.png}
	\newpage
	\thispagestyle{empty} 
	\mbox{}
	\maketitle
	\newpage
	
	\thispagestyle{empty}
	\textbf{{\LARGE ശ്രീ നാരായണഗുരു}}\\	
	\\
	തരം : ജീവചരിത്രം\\
	രചയിതാവ് : എൻ കുമാരനാശാൻ\\
	ഭാഷ : മലയാളം\\
	പ്രസിദ്ധീകരിച്ചത് : 1915\\
	പകർപ്പവകാശം : പൊതുസഞ്ചയത്തിൽ പെട്ടത്\\
	ടൈപ് ചെയ്തത് : മലയാളം വിക്കിസമൂഹം\\
	മുൻകവർ ചിത്രം : രാജശേഖരൻ (ഓയിൽ പെയ്‌ന്റിങ്) - വിക്കിമീഡിയ കോമൺസ്.\\
	വെബ് ലിങ്ക് : \href{https://ml.wikisource.org/wiki/%E0%B4%B6%E0%B5%8D%E0%B4%B0%E0%B5%80_%E0%B4%A8%E0%B4%BE%E0%B4%B0%E0%B4%BE%E0%B4%AF%E0%B4%A3_%E0%B4%97%E0%B5%81%E0%B4%B0%E0%B5%81}{ശ്രീ നാരായണ ഗുരു}
	
	\newpage
	
	\thispagestyle{empty}
	\begin{center}
		\textbf{{\LARGE പ്രസാധകക്കുറിപ്പ്}}
	\end{center}
	നാരായണഗുരുവിനെ തിരിച്ചറിഞ്ഞ ശിഷ്യപ്രധാനിയാണ് മഹാകവി കുമാരനാശാൻ. ഗുരു ജീവിച്ചിരിക്കുമ്പോൾ പ്രസിദ്ധപ്പെടുത്തിയ ഏക ജീവചരിത്രമാണിത്. 1090ൽ വിവേകോദയം മാസികയിലൂടെയാണിത് പ്രസിദ്ധപ്പെടുത്തിയത്. പിന്നീട് പ്രസിദ്ധപ്പെടുത്തിയ എല്ലാ ഗുരു ചരിത്രങ്ങൾക്കും അടിസ്ഥാനമായത് ഈ ലഘു ഗ്രന്ഥമാണ്. വിവേകോദയത്തിലൂടെ ഗുരുവിനെപ്പറ്റി ആശാൻ എഴുതിയ മുഖപ്രസംഗങ്ങളും ഈ പുസ്തകത്തിൽ ചേർത്തിട്ടുണ്ട്.
	
	
	\begin{flushright}
		വൈ.എ.റഹിം
	\end{flushright}
	
	\newpage
	\thispagestyle{empty}
	\mbox{}
	
	\newpage
	\chapter*{അദ്ധ്യായം ഒന്ന്}	
	തിരുവനന്തപുരത്തു നിന്നു മൂന്നുനാഴിക വടക്കാണ് പ്രസിദ്ധമായ ഉള്ളൂർ സുബ്രഹ്മണ്യക്ഷേത്രം. അവിടെ നിന്നു രണ്ടു നാഴിക വടക്കു കിഴക്കായി പോയാൽ ചെമ്പഴന്തി എന്ന ചരിത്രപ്രസിദ്ധമായ പഴയ ഗ്രാമമാണ്. അവിടെ ഒരു പുരാതനമായ ഈഴവ കുടുംബത്തിൽ കൊല്ലവർഷം 1032-ാമാണ്ടു ചിങ്ങമാസത്തിൽ ചതയം നക്ഷത്രത്തിൽ സ്വാമി ജനിച്ചു. മാതാപിതാക്കന്മാർ സദ്-വൃത്തിയും ഈശ്വരഭക്തിയും ഉള്ളവർ ആയിരുന്നു. അച്ഛൻ മാടനാശാൻ എന്ന ഒരു അദ്ധ്യാപകനും, അമ്മാവൻ കൃഷ്ണൻ വൈദ്യൻ എന്ന ഒരു ചികിത്സകനും ആയിരുന്നു.
	
	സ്വാമിക്കു മൂന്നു സഹോദരിമാർ ഉണ്ടായിരുന്നു. സ്വാമി കുട്ടിക്കാലത്തിൽ ശാന്തനായിരുന്നില്ല. ചൊടിപ്പുള്ള ഒരു കുട്ടിയായിരുന്നു. ചില സംഗതികളിൽ ഒരു വിധം വികൃതിയായിരുന്നു എന്നു കൂടിപ്പറയാം.
	
	വീട്ടിൽ പൂജയ്ക്കായി ഒരുക്കിവയ്ക്കുന്ന പഴവും പലഹാരങ്ങളും പൂജകഴിയുന്നതിനുമുമ്പ് എടുത്തു ഭക്ഷിച്ചു കളയുന്നതിൽ കുട്ടി അസാമാന്യമായ കൗതുകം കാണിച്ചു.
	
	`താൻ സന്തോഷിച്ചാൽ ദൈവവും സന്തോഷിക്കും' എന്നു പറയുകയും തന്റെ ആ അകൃത്യത്തെ തടയാൻ ശ്രമിക്കുന്നവരെ എങ്ങനെയെങ്കിലും ആ ബാലൻ തോല്പിക്കുകയും ചെയ്യും. തീണ്ടാൻ പാടില്ലാത്ത കീഴ്ജാതിക്കാരെ ദൂരത്തെവിടെയെങ്കിലും കണ്ടാൽ ഓടിയെത്തി അവരെ തൊട്ടിട്ട് കുളിക്കാതെ അടുക്കളയിൽ കടന്ന് സ്ത്രീകളെയും അധികം ശുദ്ധം ആചരിക്കാറുള്ള പുരുഷന്മാരെയും തൊട്ട് അശുദ്ധമാക്കുന്നതു കുട്ടിക്കു രസകരമായ ഒരു വിനോദമായിരുന്നു.
	
	ബുദ്ധിമാനും സുന്ദരനും തറവാട്ടിലെ ഏകപുത്രനുമായ കുട്ടിയെ ആ വക കുറ്റങ്ങൾക്ക് മാതാപിതാക്കന്മാർ തല്ലിയിട്ടില്ല.
	
	\newpage
	\chapter*{അദ്ധ്യായം രണ്ട്}
	
	സ്വാമിയെ വിദ്യാരംഭം ചെയ്യിച്ചത് കേൾവിപ്പെട്ട ചെമ്പഴന്തിപ്പിള്ളമാരുടെ തറവാട്ടിലെ അന്നത്തെ കാരണവരാണ്. അദ്ദേഹം ഒരു നല്ല ജ്യോത്സനും സ്ഥലത്തെ പാർവ്വത്യകാരനുമായിരുന്നു. അക്ഷരാഭ്യാസവും അന്നത്തെ ഉൾനാട്ടിലെ രീതി അനുസരിച്ച് സിദ്ധരൂപം, ബാല പ്രബോധനം, അമരം മുതലായ ബാലപാഠങ്ങളും കഴിഞ്ഞശേഷം, സ്വദേശത്ത് ഉയർന്നതരം പഠിത്തത്തിനു സൗകര്യം ഇല്ലാതിരുന്നതിനാൽ സ്വാമിക്കു പഠിപ്പു മതിയാക്കേണ്ടിവന്നു.
	
	കുടുംബത്തിലെ പ്രധാനജോലി കൃഷിയായിരുന്നു. സ്വാമി തന്റെ പ്രായത്തിനനുസരിച്ച് അതിൽ സഹായിക്കയും അടുത്ത കാട്ടുപ്രദേശങ്ങളിൽ കന്നു കാലികളെ മേച്ചുകൊണ്ടു കുറേ നാൾ കഴിക്കുകയും ചെയ്തിരുന്നു. മദ്ധ്യാഹ്നകാലത്ത് വൃക്ഷങ്ങളുടെ തണലുകളിൽ പശുക്കൾ മേഞ്ഞു നില്ക്കുമ്പോൾ സ്വാമി ഇലകൾ നിറഞ്ഞ മരക്കൊമ്പുകളിൽ കയറിയിരുന്നു നീലവർണ്ണമായ ആകാശത്തെ നോക്കി മനോരാജ്യം ചെയ്യുകയും സംസ്കൃതപദ്യങ്ങൾ ഉരുവിട്ടു പഠിക്കയും ചെയ്ക പതിവായിരുന്നു.
	
	സസ്യങ്ങൾ കൃഷിചെയ്തുണ്ടാക്കുന്നതിൽ സ്വാമിക്ക് വലിയ വാസനയായിരുന്നു. തന്നത്താൻ വെറ്റിലക്കൊടി നട്ടു നനച്ചു വളർത്തിയുട്ടുള്ളതിനെപ്പറ്റി പലപ്പോഴും സ്വാമി പറഞ്ഞു രസിക്കാറുണ്ട്. ബാല്യം കഴിയുന്നതിനു മുമ്പുതന്നെ സ്വാമി ഒരു വലിയ സാത്വികനും ഭക്തനുമാണെന്നു ജനങ്ങൾ അറിഞ്ഞു കഴിഞ്ഞു.
	
	21-ാമത്തെ വയസ്സിൽ, അതായത് 1053-ൽ, സ്വാമി സംസ്കൃതം പഠിപ്പാനായി, കരുനാഗപ്പള്ളി താലൂക്കിൽ പുതുപ്പള്ളി കുമ്മംപള്ളിൽ രാമൻപിള്ള ആശാൻ അവർകളുടെ അടുക്കലേക്കു പോയി. അവിടെ വാരണപ്പള്ളി എന്ന പ്രസിദ്ധ കുടുംബത്തിലാണ് സ്വാമി താമസിച്ചിരുന്നത്.
	
	കഴിഞ്ഞുപോയ തിരുവനന്തപുരം പെരുനെല്ലി കൃഷ്ണൻവൈദ്യർ, വെളുത്തേരി കേശവൻ വൈദ്യർ മുതലായി പലേ യോഗ്യന്മാരും സഹാധ്യായികളായിരുന്നു. ഈ സഹപാഠികളുടെയും അവിടെ ഉണ്ടായിരുന്ന മറ്റു ചെറുപ്പക്കാരുടെയും സഹവാസത്തിനേക്കാൾ സ്വാമി അധികം ഇഷ്ടപ്പെട്ടിരുന്നത് ഏകാദശി മുതലായ വ്രതങ്ങൾ അനുഷ്ഠിച്ചും പുരാണങ്ങൾ വായിച്ചും ദിവസം കഴിച്ചിരുന്ന അവിടത്തെ വൃദ്ധന്മാരുടെയും വൃദ്ധകളുടെയും സാഹചര്യത്തെ ആയിരുന്നു. ഉറക്കത്തിൽ പോലും ഈശ്വരനാമങ്ങളും മന്ത്രങ്ങളും സ്വാമിയുടെ മുഖത്തുനിന്നും സ്വതേ പുറപ്പെടുന്നതായി പലരും കേട്ടിട്ടുണ്ട്. സ്വാമിയുടെ ഈശ്വരഭക്തിയേയും സാത്വികമായ സ്വഭാവവിശേഷത്തേയും പറ്റി പല കഥകളും ആ സ്ഥലത്തുള്ളവർ ഇന്നും ഭക്തിബഹുമാനപൂർവ്വം പറഞ്ഞുവരുന്നുണ്ട്. ഗജേന്ദ്രമോക്ഷം കഥയെ വാരണപ്പള്ളി തറവാട്ടിലെ കാരണവരുടെ ആവശ്യപ്രകരം സ്വാമി ഒരു വഞ്ചിപ്പാട്ടായി അവിടെവച്ചു എഴുതിയത് ഇപ്പോഴും അവിടത്തുകാരിൽ ചിലർ പാടി കേൾക്കാറുണ്ട്. സ്വാമിയുടെ ഇഷ്ടദേവത അപ്പോൾ വിഷ്ണുവായിരുന്നു. ബാലകൃഷ്ണനെ പലപ്പോഴും മുൻപിൽ കൂത്താടുന്നതായി സ്വാമി പ്രത്യക്ഷത്തിൽ കണ്ടിട്ടുണ്ടത്രേ. സംസ്കൃതത്തിൽ പല വിഷ്ണുസ്തോത്രങ്ങളും സ്വാമി അന്ന് എഴുതിയിട്ടുണ്ടായിരുന്നു.
	
	സ്വാമിയുടെ ബുദ്ധിയും ഓർമ്മയും വലിയ ശക്തിയുള്ളവയായിരുന്നു. ഒരിക്കൽ വായിച്ച പുസ്തകങ്ങളേയോ കേട്ട വിഷയത്തേയോ മറക്കുക പതിവില്ലായിരുന്നു. രണ്ടുകൊല്ലത്തിനുള്ളിൽ കാവ്യനാടകാലംകാരങ്ങളിൽ നല്ല വ്യുത്പത്തി സമ്പാദിച്ച് ഗുരു ദക്ഷിണകഴിഞ്ഞ് അവിടത്തെ പഠിത്തം അവസാനിപ്പിച്ചു. ഗുരുനാഥനായ രാമൻപിള്ള ആശാൻ അവർകൾക്ക് എല്ലാ ശിഷ്യന്മാരിലും വച്ച് സാത്വികനായസ്വാമിയിൽ സ്നേഹവിശേഷം ഉണ്ടായിരുന്നു. വാരണപ്പള്ളിയിൽ നിന്നും മടങ്ങിപ്പോരാൻ ഒരുങ്ങുമ്പോൾ ഒരു കഠിനമായ രക്താതിസാരം ആരംഭിച്ച്തിനാൽ സ്വദേശത്തിൽ നിന്ന് ആളുകൾചെന്ന് സ്വാമിയെ കൊണ്ടുപോരുകയാണുണ്ടായത്. സ്ഥലം വിടുമ്പോൾ രോഗത്തിന്റെ കാഠിന്യത്താൽ സ്വാമിക്കു പ്രജ്ഞയില്ലായിരുന്നു. സ്വാമിയുടെ ആ സ്ഥിതിയിലുള്ള വേർപാടിൽ കരയാത്തവരായി അന്നു അവിടെ ആരും ഉണ്ടായിരുന്നില്ല.
	
	\newpage
	
	\chapter*{അദ്ധ്യായം മൂന്ന്}
	
	വീട്ടിൽ മടങ്ങിവന്നതിനുശേഷം അധികം താമസിയാതെ തന്നെ രോഗം സുഖപ്പെട്ടു. കുറേക്കാലം ചെമ്പഴന്തിയിൽ തന്നെ സ്വാമി കുട്ടികളെ വായിപ്പിച്ചു താമസിച്ചു. നാണു ആശാൻ എന്ന പഴയ പേർ സ്വാമിക്ക് അങ്ങനെ സിദ്ധിച്ചതാണ്. നാൾ പോകും തോറും സ്വാമിക്ക് ഈശ്വരഭക്തി വർദ്ധിക്കുകയും ലൗകിക ജീവിതത്തിൽ ശക്തി കുറഞ്ഞുവരികയും ചെയ്തു. "ഗീതാഗോവിന്ദം" എന്ന പ്രസിദ്ധഗ്രന്ഥം അക്കാലത്തു സ്വാമി ദിവസേന പാരായണം ചെയ്തിരുന്നതായി അറിയുന്നു. ഇതിനിടയിൽ മാതാപിതാകന്മാരുടേയും മറ്റും നിർബന്ധത്താൽ സ്വാമിക്ക് വിവാഹം കഴിക്കേണ്ടി വന്നു. എന്നാൽ വിഷയസുഖങ്ങൾക്കു സ്വാമിയെ വ്യാമോഹിപ്പിപ്പാൻ കഴിഞ്ഞില്ല. പിന്നെ സ്വല്പ കാലത്തിനുള്ളിൽ അമ്മ മരിച്ചു. 1060-ൽ അച്ഛനും കാലധർമ്മം പ്രാപിച്ചു. സ്വാമി ദാമ്പത്യ ബന്ധത്തിൽ നിന്ന് ഇതിനുമുമ്പുതന്നെ മോചിച്ചിരുന്നു. മാതാപിതാക്കന്മാരുടെ മരണശേഷം കുടുംബ ബന്ധത്തിൽ നിന്നും മോചിക്കുന്നത് സ്വാമിക്ക് എളുപ്പമായിത്തീർന്നു. കാരണവർ ഗൃഹഭരണത്തിനായി സ്വാമിയെ നിർബന്ധിച്ചു എങ്കിലും മതസംബന്ധമായ വിഷയത്തിൽ ജീവിതകാലം നയിപ്പാനായുള്ള തന്റെ ദൃഡനിശ്ചയത്തിന് ആ നിർബന്ധംകൊണ്ട് ഇളക്കം ഒന്നും ഉണ്ടായില്ല. സ്വാമി ഇതിനുശേഷം വീട്ടിൽ താമസിക്കാതെയായി. രാത്രിയും പകലും സമീപത്തുള്ള ജനവാസമില്ലാത്ത കുന്നുകളിലും, കാടുകളിലും, പാറ ഇടുക്കുകളിലും, സമുദ്രതീരങ്ങളിലും, ക്ഷേത്രങ്ങളിലും, മറ്റ് ഏകാന്തസ്ഥലങ്ങളിലും സ്വാമി ധ്യാനനിരതനായി ഇരിക്കുന്നതും, എകാകിയായി സഞ്ചരിക്കുന്നതും പലപ്പോഴും പലരും കണ്ടിട്ടുണ്ട്.
	
	ഇക്കാലത്ത് സ്വാമി, 'പ്രാചിന മലയാളം' മുതലായ ഗ്രന്ഥങ്ങലുടെ കർത്താവായ കുഞ്ഞൻപിള്ള ചട്ടമ്പി എന്ന മഹാനുമായി പരിചയപ്പെടുകയും ആ വഴി തിരുവനന്തപുരത്തു "തൈക്കാട്ട് അയ്യാവ്" എന്ന സുബ്രഹ്മണ്യഭക്തനും യോഗിയുമായ ഗുരുവിന്റെ അടുക്കൽ നിന്ന് യോഗാഭ്യാസസംബന്ധമായ ഉപദേശം കൈക്കൊള്ളുകയും ചെയ്തു. സ്വാമി സുബ്രഹ്മണ്യോപാസകനായത് ഇതുമുതലാണ്. ഇങ്ങനെ രണ്ടുമൂന്നു കൊല്ലം സ്വദേശത്തുതന്നെ പല സ്ഥലങ്ങളിലുമായി യോഗം ശീലിച്ചുകൊണ്ട് താമസിക്കയും സഞ്ചരിക്കയും ചെയ്തു. വേളിയിൽ സമുദ്രതീരത്ത് ഒരു കുടിൽ കെട്ടി സ്വാമി കുറെനാൾ അതിൽ താമസിക്കയും അതിനുശേഷം അഞ്ചുതെങ്ങിൽ ഒരു ഒഴിഞ്ഞ പഴയക്ഷേത്രത്തിൽ കുറേനാൾ ഇരിക്കുകയും ചെയ്തിട്ടുണ്ട്. അവിടെ ചിലരെ ഇഷ്ടമുള്ള സമയങ്ങളിൽ സ്വാമി സംസ്കൃതം വായിപ്പിച്ചിരുന്നു. സ്വാമിയെ അക്കാലങ്ങളിൽ ഏതാനും ദിവസം സ്ഥിരമായി ഒരു ദിക്കിൽ കാണുക പതിവില്ല. കുറേനാൾ സ്വദേശത്ത് എങ്ങും തന്നെ കണ്ടില്ല. അപ്പോൾ ദക്ഷിണ ഇന്ത്യയിലുള്ള പഴനി തുടങ്ങിയ പല മഹാക്ഷേത്രങ്ങളും സ്വാമി സന്ദർശിച്ചിട്ടുള്ളതായറിയുന്നു. ആ സഞ്ചാരത്തിൽ സ്വാമി ഭിക്ഷാന്നംകൊണ്ടു ഉപജീവിക്കയും, വഴിയമ്പലങ്ങളിൽ കിടന്നുറങ്ങുകയും ആണ് ചെയ്തിട്ടുള്ളത്. ഇങ്ങനെയുള്ള സഞ്ചാരങ്ങളിൽ സ്വാമിക്കു പലപ്പോഴും ആപത്തുകൾ നേരിടാൻ പോയതായും, അതിൽനിന്നൊക്കെയും അൽഭുതകരമാം വണ്ണം രക്ഷപ്പെട്ടിട്ടുള്ളതായും പല കഥകളും കേട്ടിട്ടുണ്ട്.
	
	\newpage
	
	\chapter*{അദ്ധ്യായം നാല്}
	
	കുറെക്കാലം കഴിഞ്ഞു സ്വാമി മടങ്ങിയെത്തി. എകാകിയായും അജ്ഞാതനായും കേരളത്തിന്റെ നാനാഭാഗത്തും സ്വാമി ഇക്കാലത്തു സഞ്ചരിച്ചിരുന്നു. സ്വാമിയുടെ ദൃഢമായ ബ്രഹ്മചര്യവും, തപസ്സും, യോഗവും ക്രമേണ അൽഭുതകരങ്ങളായ ഫലങ്ങളെ പ്രദർശിപ്പിച്ചുതുടങ്ങി. ചെല്ലുന്ന ദിക്കിലെല്ലാം കുഷ്ഠം മുതലായ മഹാരോഗങ്ങൾ പിടിപ്പെട്ടവർ സ്വാമിയുടെ അടുക്കൽ വന്നുചേരുകയും, ഏതെങ്കിലും ഒരു പച്ചിലയോ എന്തെങ്കിലും ഒരു ഭക്ഷണസാധനമോ എടുത്തുകൊടുത്ത് അവരുടെ രോഗങ്ങളെ സ്വാമി സുഖപ്പെടുത്തുകയും ചെയ്തിരുന്നു. ബ്രഹ്മരക്ഷസ്സ്, അപസ്മാരം മുതലായ ഉപദ്രവത്താൽ വളരെക്കാലം കഷ്ടത അനുഭവിച്ചിരുന്ന രോഗികൾക്കു സ്വാമിയുടെ ദർശനമാത്രത്താൽതന്നെ പൂർണ്ണസുഖം കിട്ടിയിട്ടുണ്ട്. അനവധി കുട്ടിച്ചാത്തന്മാരുടെ ഉപദ്രവങ്ങളെയും സ്വാമി വിലക്കി മാറ്റിയിട്ടുണ്ട്. നാനാജാതിക്കാരായ ഹിന്ദുക്കളുടെ ഇടയിൽ പല മാന്യതറവാടുകളിലുമുള്ള വന്ധ്യകൾ സ്വാമിയുടെ കൈകൊണ്ടു വല്ല പഴമോ മറ്റോ വാങ്ങി ഭക്ഷിക്കുകയോ സ്വാമിയുടെ ഒരു അനുഗ്രഹവാക്കു ലഭിക്കുകയോ ചെയ്തശേഷം താമസിയാതെ ഗർഭംധരിച്ചു പ്രസവിച്ചിട്ടുണ്ട്. ഒരിക്കലും കുടിവിടാത്ത അനേകം മദ്യപന്മാരുടെ മദ്യപാനവും സ്വാമി നിർത്തിയിട്ടുണ്ട്. സ്വാമിയുടെ വാക്കിനെ ലംഘിച്ചു കൊതികൊണ്ടു വീണ്ടും മദ്യപാനം ആരംഭിച്ച ചിലർ മദ്യം കാണുമ്പോൾ ഛർദ്ദിക്കുകയും ചെയ്തിട്ടുണ്ട്.
	
	സ്വാമിയുടെ ഈ അദ്ഭുതപ്രവൃത്തികളും സാത്വികനിഷ്ഠയും കണ്ടു വിശ്വാസത്താൽ പലരും അദ്ദേഹത്തിന്റെ പേരിൽ ഈശ്വരന്റെ പേരിൽ എന്ന പേലെ നേർച്ചകൾ നേരുകയും രോഗശാന്തി മുതലായ ഫലങ്ങൾ അതിൽനിന്നും അവർക്കു ലഭിക്കുകയും ചെയ്തിട്ടുണ്ട്. ആ താരുണ്യദശയിൽ പ്രകൃത്യാ സുമുഖനും ശാന്തഹൃദയനുമായ സ്വാമി ഏതു ജനക്കൂട്ടത്തിനിടയിൽ കാണപ്പെട്ടാലും യോഗശക്തികൊണ്ട് ഉജ്ജലമായ മുഖത്തുള്ള പ്രത്യേക തേജോവിശേഷം അവിടന്ന് ഒരു അമാനുഷനാണെന്നു വിളിച്ചു പറയുമായിരുന്നു. സ്വാമി ജനങ്ങളുടെ ഇടയിൽ നിന്നും തെറ്റിഒഴിഞ്ഞ് എകാന്തമായി സഞ്ചരിക്കുക സാധാരണയായിരുന്നു. അക്കാലത്തു ഗ്രാമങ്ങളിലോ നഗരങ്ങളിലോ അദ്ദേഹത്തെ കണ്ടെത്തിയാൽ ഒരു വലിയ പുരുഷാരം ചുറ്റും കൂടുക പതിവാണ്. സ്വാമി അക്കാലത്ത് സംസ്കൃതത്തിലും മലയാളത്തിലും അതിമനോഹരങ്ങളായ പല സുബ്രഹ്മണ്യസ്തോത്രങ്ങളും എഴുതിയിട്ടുണ്ട്. അന്നു പശ്ചാത്യ വിദ്യാഭ്യാസത്തിന്റെ പുതിയ ആവിർഭാവത്തോടുകൂടി നാട്ടിൽ ബാധിച്ചിരുന്ന നാസ്തിക വിചാരങ്ങൾക്കു സ്വാമിയുടെ ജീവിതം തന്നെ പലർക്കും ഒരു പരിഹാരമായിരുന്നു എന്നുള്ളതും പ്രസ്താവയോഗ്യമാണ്.
	
	ഈയിടയിൽ കുറെക്കാലം സ്വാമി ഭക്ഷ്യപേയങ്ങളെ സംബന്ധിച്ചു ജാതിഭേദമോ വകഭേദമേ വിചാരിക്കാതെ കൊടുക്കുന്നവരുടെ കയ്യിൽനിന്നും കിട്ടുന്നതെല്ലാം വാങ്ങിക്കഴിച്ചിരുന്നു. ഇങ്ങനെ ചില വിഷഭക്ഷണങ്ങൾ പോലും സ്വാമി ഭക്ഷിച്ചിട്ടു യാതൊരു വ്യാപത്തും ഉണ്ടാകാതെ ഇരുന്നതായി കേട്ടിട്ടുണ്ട്.
	
	\newpage
	
	\chapter*{അദ്ധ്യായം അഞ്ച്}
	
	അങ്ങനെ ഇരിക്കുമ്പോഴാണ് 1063-ാമാണ്ടിടയ്ക്ക് ഒരു കാട്ടുപ്രദേശമായിരുന്ന അരുവിപ്പുറം സ്വാമി സന്ദർശിച്ചത്. ജനവാസമില്ലാത്ത ആ സ്ഥലത്തെ മനോഹരമായ നദീപ്രവാഹം പാറകളിൽ തടഞ്ഞുണ്ടാകുന്ന ഗംഭീരമായ മുഴക്കവും, കരയിലെ പാറ ഇടുക്കുകളും, മണൽതിട്ടകളും, രണ്ടു വശത്തുള്ള ഉന്നതമായ കുന്നുകളും വൃക്ഷലതാതികൾ നിറഞ്ഞ പച്ചനിറമായ കാടുകളും ഏകാന്തപ്രിയനായ സ്വാമിയെ സാമാന്യത്തിലധികം ആകർഷിച്ചു. സ്വാമി ചില അവസരങ്ങളിൽ അവിടെയുള്ള പാറയിടുക്കുകളിൽ അനേകദിവസം തുടർന്നുകൊണ്ടു യാതൊരു ആഹാരവും കൂടാതെയും താൻ അവിടെയുള്ള വിവരം ആരും അറിയാതെയും കഴിച്ചുകൂട്ടിയിട്ടുണ്ട്. ക്രമേണ ജനങ്ങൾ സ്വാമിയുടെ അവിടെയുള്ള സഞ്ചാരത്തേയും സങ്കേതസ്ഥലങ്ങളേയും മനസിലാക്കി, ചില ഭക്തന്മാർ അടുത്ത ഗ്രാമത്തിൽ നിന്നു ചിലപ്പോഴെല്ലാം ഭക്ഷണങ്ങൾ കൊണ്ടുവന്നു കൊടുത്തുതുടങ്ങി. ഇതിനുശേഷം അരുവിപ്പുറം അധികകാലം ഒരു ഏകാന്ത സ്ഥലമായിരുന്നില്ല. പല ദിക്കുകളിൽ നിന്നും സ്വാമിയെ അന്വേഷിച്ച് ആളുകൾ അവിടെ എത്തിത്തുടങ്ങി. രോഗം ശമിപ്പിക്കുകയും ഭൂതങ്ങളെ ഒഴിക്കുകയും ഉപദേശം നൽകുകയും ശാസ്ത്രാർഥങ്ങൾ പറഞ്ഞു കൊടുക്കുകയും മറ്റു പ്രകാരത്തിൽ സാധുക്കളെ അനുഗ്രഹിക്കുകയും ചെയ്യേണ്ട ഭാരം ആ വിജനത്തിലും സ്വാമിയെ പിൻതുടർന്നു. സ്വാമിയുടെ ദൂരെദർശനം, പരഹൃദയജ്ഞാനം മുതലായ സിദ്ധികളുടെ പല ദൃഷ്ടാന്തങ്ങൾ അക്കാലത്തു ജനങ്ങൾ അനുഭവിച്ചറിഞ്ഞിട്ടുള്ളതാണ്. നാനാജാതിക്കാരായ പലഭക്തന്മാർ സ്വാമിയെ സന്ദർശിപ്പാൻ അവിടെ വരികയും അവരിൽ ചിലർ സ്വാമിയുടെ ശിഷ്യന്മാരായി തീരുകയും ചെയ്തു. പല ദിക്കുകളിലുള്ള ഭക്തന്മാരായ ഗൃഹസ്ഥന്മാർ അരി മുതലായ സംഭാരങ്ങളോടുകൂടി അവിടെവന്നു സ്വാമിക്കും കൂടെയുള്ളവർക്കും സദ്യകൾ കഴിച്ചു മടങ്ങിപ്പോവുക പതിവായിരുന്നു. ഇങ്ങനെ സ്വാമിയുടെ സാന്നിദ്ധ്യം കൊണ്ടുതന്നെ അവിടം വേഗത്തിൽ ഒരു പുണ്യസ്ഥലമായി സ്ഥലമായി പരിണമിച്ചു. സ്വാമി ഇല്ലാത്ത സമയങ്ങളിൽ കൂടിയും ജനങ്ങൾ അവിടെ വന്നു സ്നാനം ചെയ്തു തൊഴുതുപോകുവാൻ തുടങ്ങി. അതിനുശേഷം അവിടെ ഒരു ആരാധനാ സ്ഥലം ഉണ്ടായിരുന്നാൽ കൊള്ളാമെന്നു സ്വാമി ചിലരോടു പറയുകയും അങ്ങനെ എന്തെങ്കിലും ചെയ്യുന്നതിൽ തനിക്കുള്ള ആഭിമുഖ്യത്തെ പതിവായി തന്നെ വന്നുകാണുന്ന ഭക്തന്മാരായ ചില ചെറുപ്പക്കാരോടു സൂചിപ്പിക്കുകയും ചെയ്തു. 1063-ാമാണ്ടത്തെ ശിവരാത്രി സമീപിച്ചാണ് സ്വാമിയുടെ ഈ അഭിലാഷം പുറത്തായത്. വിഗ്രഹങ്ങൾ ഉണ്ടാക്കുവാനോ ക്ഷേത്രങ്ങളോ കെട്ടിടങ്ങളോ പണിയിക്കാനോ ആ കാട്ടുപ്രദേശത്ത് യാതൊരു സൗകര്യവും ഉണ്ടായിരുന്നില്ല. സ്വാമി അതൊന്നും ആവിശ്യപ്പെട്ടുമില്ല. നദിയുടെ കിഴക്കെ തീരത്തുള്ള ഒരു പാറയെ പീഠമായി സങ്കല്പിച്ച് അതിന്മേൽ ഏതാണ്ട് ഒരു ശിവലിംഗാകൃതിയിൽ ആറ്റിൽ കിടന്ന ഒരു ശിലാഖണ്ഡം എടുത്ത് ശിവരാത്രി ദിവസം പ്രതിഷ്ഠ നടത്താനാണ് സ്വാമിയുടെ ഭാവം എന്നറിയുകയും അടുത്തുള്ള ജനങ്ങൾ തങ്ങളാൽ കഴിയുന്ന ചില ചെറിയ ഒരുക്കങ്ങൾ അതിനായി ചെയ്യുകയും ചെയ്തു. സ്വാമി ഇരിക്കുന്നതായറിഞ്ഞ് ഏതാനും വ്രതക്കാരും ഭജനക്കാരും ശിവരാത്രിനാൾ ``ഉറക്കിളയ്ക്കാ"നായി അവിടെകൂടി. പ്രതിഷ്ഠയുടെ സംഭാരമായി ഉണ്ടായിരുന്നത് കുറെ പുഷ്പങ്ങളും വിളക്കുകളും നാദസ്വരവായനയും മാത്രമായിരുന്നു. പീഠമായി സങ്കൽപ്പിച്ചിരുന്ന പാറയുടെ മീതെ ഒരു ചെറിയ പന്തൽ കെട്ടിയിരുന്നു. അർദ്ധരാത്രിയോടുകൂടി സ്വാമി സ്നാനം ചെയ്തുവന്ന് അതിനകത്തു കടന്നു. പ്രതിഷ്ഠിക്കാനുള്ള ശിലയെ കൈയിൽ എടുത്തു ധ്യാനിച്ചു കൊണ്ടു രാത്രി മൂന്നുമണിവരെ ഒരേ നിലയിൽ നിന്നു. സ്വാമിയുടെ തേജോമയമായ മുഖത്ത് ആ സമയം അശ്രുധാരകൾ പ്രവഹിച്ചുകൊണ്ടിരുന്നു. കാണികൾ ഭക്തിപരവശന്മാരായി പഞ്ചാക്ഷരമന്ത്രം ഉച്ചത്തിൽ ജപിച്ച് ഏക മനസ്സോടെ ചുറ്റും നിന്നു. മൂന്നുമണിക്ക് ശിലയെ സ്വാമി പീഠത്തിൽ പ്രതിഷ്ഠിച്ച് അഭിഷേകം ചെയ്തു. ആ സമയത്തു ചില അൽഭുതങ്ങൾ കണ്ടിട്ടുള്ളതായി പലരും പറയുന്നു. ഇങ്ങനെയാണ് സ്വാമിയുടെ മതസംബന്ധമായ സ്ഥാപനങ്ങളിൽ ആദ്യത്തേതായ അരുവിപ്പുറം ക്ഷേത്രം ആരംഭിച്ചത്.
	
	ഈ പ്രതിഷ്ഠയ്ക്കുശേഷം സ്വാമി മുൻപിലത്തേക്കാൾ അധികം അരുവിപ്പുറത്ത് താമസിക്കുക പതിവായി. ക്ഷേത്രത്തിന്റെ സന്നിദ്ധ്യം ക്രമേണ വർദ്ധിച്ചു. അവിടെ വരുന്ന കാണിക്കകൾ എടുത്തു മുതൽക്കൂട്ടിയും സംഭാവനകൾ പിരിച്ചും ജനങ്ങൾ പ്രതിഷ്ഠാസ്ഥാലത്ത് ഒരു ശ്രീകോവിൽ കെട്ടിച്ചു. ചില ഭക്തന്മാരായ ശിഷ്യന്മാരെ സ്വാമി ഷേത്രത്തിൽ ശാന്തിക്കാരാക്കി. പഠിപ്പുള്ള ശിഷ്യന്മാരിൽ ചിലരെക്കൊണ്ടു അവിടെ ഒരു പള്ളിക്കൂടം കെട്ടിച്ച് അടുത്തുള്ള കുട്ടികളെ മലയാളം പഠിപ്പിക്കാൻ ഏർപ്പാടുചെയ്തു. ഇക്കാലത്താണ് സ്വാമി പ്രസിദ്ധമായ 'ശിവശതകം' എന്ന മണിപ്രവാളസ്തോത്രം നിർമ്മിച്ചത്.
	
	\newpage
	
	\chapter*{അദ്ധ്യായം ആറ്}
	
	ഇതിനു മുൻപേ തന്നെ സ്വാമിക്ക് അദ്വൈതശസ്ത്രത്തിൽ പ്രതിപത്തി വർദ്ധിച്ചിരുന്നു. പല വേദാന്തികളുമായി സ്വാമി സഹവാസം ചെയ്കയും ചെയ്തിരുന്നു. കന്യാകുമാരിക്ക് സമീപമുള്ള ശ്രുതിപ്പെട്ട മരുത്വാമലയിൽ സ്വാമി കൂടക്കൂടെ പോയിവന്നിരുന്നു. നാഗരുകോവിലിനു സമീപം ഒരിടത്ത് അപ്പോൾ ഒരു യോഗിനി നിർവ്വികൽപ്പസമാധിയിൽ കിടന്നിരുന്നത് സ്വാമി പലപ്പോഴും സന്ദർശിച്ചിട്ടുണ്ട്. ആ മഹാനുഭാവ സ്വാമിയെ ഒരിക്കൽ അനുഗ്രഹിച്ചിട്ടുള്ളതായറിയുന്നു.
	
	ഇക്കാലത്തു സ്വാമി തമിഴിൽ നല്ല അറിവു സമ്പാദിക്കയും വേദാന്ത വിഷയമായി അനേകം പാട്ടുകൾ ആ ഭാഷയിൽ എഴുതുകയും ചെയ്തിട്ടുണ്ട്. ``തിരുക്കുറളും" ``ഒഴിവിലൊടുക്ക"മെന്ന പ്രൗഢഗ്രന്ഥവും സ്വാമി തമിഴിൽ നിന്നും മലയാളത്തിൽ പദ്യരൂപേണ ഭാഷാന്തരപ്പെടുത്തിയിട്ടുള്ളതായും അറിയാം. സ്വാമിയെപ്പോലെ അസാമാന്യമായ `ബുദ്ധിശക്തിയും' ശാസ്ത്രപാണ്ഡിത്യവും, അനുഷ്ഠാനവും അനുഭവവും ഒത്തു ചേർന്ന ഒരു വേദാന്തി ദുർലഭമായിരുന്നതിനാൽ അദ്ദേഹത്തെ കാണുന്നത് വിദ്വാന്മാർക്കു കൗതുകകരമായിരുന്നു. പ്രസിദ്ധ തത്വജ്ഞാനിയായി കഴിഞ്ഞുപോയ പ്രൊഫസർ സുന്ദരംപിള്ള (എം. എ) അവർകൾ തുടങ്ങിയ യോഗ്യന്മാരുടെ അസാമാന്യമായ ശ്ലാഘയ്ക്കും ഭക്തിബഹുമാനങ്ങൾക്കും സ്വാമി പാത്രമായിത്തീർന്നിരുന്നു.
	
	\newpage
	
	\chapter*{അദ്ധ്യായം ഏഴ്}
	
	അരുവിപ്പുറത്തെ പ്രതിഷ്ഠയ്ക്കുശേഷം സ്വാമി ചിറയിൻകീഴ് വക്കത്ത് വേലായുധൻ കോവിൽ എന്ന പഴയ ഒരു സുബ്രഹ്മണ്യക്ഷേത്രം പുതുക്കി പ്രതിഷ്ഠിക്കയും വക്കത്തു തന്നെ ദേവേശ്വരം എന്ന വേറൊരു ശിവക്ഷേത്രം കൂടി പ്രതിഷ്ഠിക്കയും ചെയ്തു. സ്വജാതിക്കാരുടെ ഇടയിൽ തിരുവിതാംകൂറിന്റെ തെക്കേ അറ്റം മുതൽ വടക്കേ അറ്റം വരെ സ്വാമിയുടെ പേരു ശക്തിയോടുകൂടി പരക്കുകയും പല ഇടത്തുമുള്ള പ്രധാന യോഗ്യന്മാരെല്ലാം സ്വാമിക്കു പരിചിതരായോ ശിഷ്യന്മാരായോ തീരുകയും ചെയ്തു, അക്കാലത്തു വടക്കൻ പറവൂരിലും ആലുവായിലും സ്വാമി ഒന്നിലധികം പ്രാവിശ്യം പോയി താമസിച്ചിരുന്നു. ഇപ്പോൾ അവിടെ അദ്വൈതാശ്രമം സ്ഥാപിച്ചിരുക്കുന്ന സ്ഥലത്ത് ഒരു ആശ്രമം സ്ഥാപിച്ചാൽ കൊള്ളാമെന്ന് ആലുവാപ്പുഴയിൽ അന്നു കുളിച്ചു നിൽക്കുമ്പോൾ വിചാരിച്ചിട്ടുള്ളതായി ഒടുവിൽ സ്വാമി പ്രസ്താവിച്ചിട്ടുണ്ട്. അരുവിപ്പുറത്തു ക്ഷേത്രത്തോടു ചേർത്ത് ഒരു സന്യാസിമഠം ഉറപ്പിക്കേണമെന്നും അതുമൂലമായി ജനങ്ങളുടെ ഇടയിൽ മതസംബന്ധമായ അറിവുവർദ്ധിപ്പിച്ചു ക്ഷേത്രങ്ങളിൽ ഹിംസ മുതലായ അകൃത്യങ്ങളെ തടുക്കുകയും ദുർദേവതാരാധനകളെ നിർത്തൽ ചെയ്യുകയും സാത്വികമായ ആരാധനാക്രമങ്ങളെ പ്രചാരപ്പെടുത്തുകയും ചെയ്യണമെന്നും സ്വാമി തീർച്ചയാക്കിയതായി അവിടത്തെ പ്രവർത്തികളിൽ നിന്നു പ്രത്യക്ഷപ്പെട്ടുതുടങ്ങി. പറവൂർ മുതൽ നെയ്യാറ്റുംകരവരെ ഈഴവരുടെ പഴയ ദേവിക്ഷേത്രങ്ങളിൽ പലതിലും നടന്നിരുന്ന ആട്, കോഴി മുതലായ ജന്തുക്കളെ ബലികൊടുക്കുന്ന അനാചാരം സ്വാമിയുടെ വാക്കാൽ അന്നുമുതൽ ജനങ്ങൾ പലസ്ഥലത്തും വിട്ടുകളഞ്ഞിരുന്നു. ഹിംസയുടെയും ദുർദേവതാരാധനയുടെയും ദോഷത്തെപ്പറ്റി സ്വാമി പല സാരമായ ലഘുലേഖനങ്ങളും സ്വന്തമായി എഴുതി പ്രാസംഗികന്മാരെ ഏൽപ്പിച്ചു പരസ്യമായും മതപരിഷ്കരണവിഷയത്തിൽ പ്രസംഗങ്ങൾ നടത്തിച്ചിരുന്നു. അരുവിപ്പുറം ക്രമേണ ഒരു സന്യാസി മഠമായിത്തീർന്നു. പഠിപ്പും ഭക്തിയുമുള്ള പല ചെറുപ്പക്കാരും അവിടെവന്നു സ്വാമിയുടെ ശിഷ്യന്മാരായി താമസം തുടങ്ങി. മഠത്തെ ഒരു സ്ഥിരമായ സ്ഥാപനമാക്കേണ്ട ആവിശ്യകത വർദ്ധിച്ചുവന്നു.
	
	1068-ൽ സ്വാമി തിരുവനന്തപുരത്തിനടുത്തുള്ള കുളത്തൂർ ഈഴവരുടെ വകയും വളരെ പുരാതനവും ആയ കോലത്തുകര ഭഗവതിക്ഷേത്രം പൊളിച്ചുമാറ്റി തൽസ്ഥാനത്ത് ഒരു ശിവക്ഷേത്രം നിർമിച്ചു പ്രതിഷ്ഠകഴിച്ചു.
	
	1069-ൽ സ്വാമി അരുവിപ്പുറം മഠത്തിലെ ധർമ്മകാര്യങ്ങൾ അന്വേഷിപ്പാനും ജനങ്ങളിൽ നിന്നും പണം യാചിച്ചും മറ്റും അതിനെ വർദ്ധിപ്പിപ്പാനുമായി ചില വ്യവസ്ഥകൾ ചെയ്തു. 1070-ൽ സ്വാമി ശിഷ്യസമേതനായി മഠത്തെ അഭിവൃദ്ധിപ്പെടുത്തുന്നതിനുള്ള ചില ആവിശ്യങ്ങൾക്കായി പറവൂർവരെ സഞ്ചരിക്കുകയും അവിടെനിന്ന് ബാംഗ്ലൂർ വരെ പോയി ഡോക്ടർ പൽപ്പു അവർകളെ കണ്ടു മടക്കത്തിൽ ചിദംബരം, മധുര മുതലായ സ്ഥലങ്ങൾ സന്ദർശിച്ചു തിരുനൽവേലിവഴിയായി തിരിച്ചെത്തുകയും ചെയ്തു. ഈ അവസരത്തിൽ ശിഷ്യന്മാരിൽ ചിലരുടെ ഉയർന്ന തരം പഠിത്തങ്ങൾക്കായി സ്വാമി ഏർപ്പാടുകൾ ചെയ്കയും അവരിൽ നിന്നും മഠത്തിന്റെ അഭിവൃദ്ധിയെ പ്രതീക്ഷിക്കയും ചെയ്തിരുന്നു.
	
	തിരുവനന്തപുരത്തിന് തെക്ക് കോവളം എന്ന സുഖവാസസ്ഥലത്തിന് സമീപം മുട്ടയ്ക്കാട് എന്ന സ്ഥലത്ത് ഒരു മനോഹരമായ കുന്നിന്മേൽ സ്വാമി ഇതിനുമുമ്പുതന്നെ ഒരു സുബ്രഹ്മണ്യക്ഷേത്രത്തിന്റെ പ്രതിഷ്ഠനടത്തി. ആ ക്ഷേത്രവും അതിന്റെ ചിലവിലേക്കായി കുറെ ഭൂസ്വത്തുക്കളും അതിന്റെ ഉടമസ്ഥനായ കൊച്ചുകുട്ടി വൈദ്യർ അവർകൾ 1074-ൽ സ്വാമിക്കു ദാനം എഴുതിക്കൊടുത്തു.
	
	അവിടെ പാറയിൽ നിന്ന് നിർമ്മലജലം ഊറിക്കൊണ്ടിരിക്കുന്ന ഒരു ഉറവയുണ്ട്. അതിനടുത്ത് നല്ല ഭംഗിയും സൗകര്യവും ഉള്ള ഒരുമഠം കൂടി ടി വൈദ്യർ ഇപ്പോൾ കെട്ടിച്ചിരിക്കുന്നു.
	
	1074-ാമാണ്ട് അരുവിപ്പുറം ക്ഷേത്രത്തിന്റെയും മഠത്തിന്റേയും അഭിവൃദ്ധിയേയും ഭരണത്തേയും ഉദ്ധേശിച്ചു നെയ്യാറ്റിങ്കര, തിരുവനന്തപുരം ഈ താലൂക്കുകളിലുള്ള സ്വജനങ്ങളിൽ എതാനും മാന്യന്മാരെ കൂട്ടിച്ചേർത്ത് ``അരുവിപ്പുറം ക്ഷേത്രയോഗം" എന്ന പേരിൽ ഒരു സംഘം ഏർപ്പെടുത്തി ക്ഷേത്രത്തിലെ ഉത്സവാദികാര്യങ്ങൾ ആ സംഘംമുഖേന നടത്തിവന്നു. ഇതിനിടയ്ക്ക് മദ്ധ്യതിരുവിതാംകൂറിലും ഉത്തര തിരുവിതാംകൂറിലും ഉള്ള ചില സ്ഥലങ്ങളിൽകൂടി അതാതു സ്ഥലങ്ങളിലെ സ്വജനങ്ങളുടെ അപേക്ഷപ്രകാരം സ്വാമി ക്ഷേത്രങ്ങൾ പ്രതിഷ്ഠിച്ചു കൊടുക്കുകയും അവിടെ എല്ലാം പരിഷ്കൃതമായ ആരാധനാക്രമങ്ങൾ ഏർപ്പെടുത്തുകയും ചെയ്തു. 1076(1901)ലെ തിരുവിതാംകൂർ സെൻസസ് റിപ്പോർട്ടിൽ ഈഴവസമുദായത്തെപ്പറ്റി വിവരിക്കുന്ന ദിക്കിൽ സ്വാമിയെപ്പറ്റി "A pious religious reformer" (ഒരു സ്വാത്വികനായ മതപരിഷ്കാരൻ) എന്നു പറഞ്ഞിരിക്കുന്നതിൽ നിന്ന് സ്വാമിയുടെ മതസംബന്ധമായ പ്രവർത്തികൾക്ക് അന്നുതന്നെ ഉണ്ടായിരുന്ന പ്രസിദ്ധി ഊഹിക്കാവുന്നതാണ്.
	
	\newpage
	
	\chapter*{അദ്ധ്യായം ഏട്ട്}
	
	1078 ധനു 23-ാം തിയതി സ്വാമി തന്റെ മതസംബന്ധമായും സമുദായസംബന്ധമായും ഉള്ള ഉദ്ദേശങ്ങളെ നടപ്പിൽ വരുത്തുന്നതിനായി ``ശ്രീ നാരായണ ധർമ്മ പരിപാലനയോഗം" സ്ഥാപിക്കുകയും അരുവിപ്പുറം ക്ഷേത്രയോഗം അതിൽ ലയിക്കുകയും ചെയ്തു. ഉടനെ സ്വാമി വടക്കൻ പറവൂർ ശ്രീ നാരായണ മംഗലം ക്ഷേത്രത്തിലെ പ്രതിഷ്ഠാകർമ്മത്തിനു പോവുകയും അവിടെനിന്നു തൃപ്രയാറിനു സമീപമുള്ള പെരിങ്ങോട്ടുകര, വലപ്പാട്, വടാനപ്പള്ളി മുതലായ സ്ഥലങ്ങളിൽ പോയി വീണ്ടും പറവൂർ എത്തുകയും മകരം ഇരുപതാം തിയതിയോടു കൂടി പ്രതിഷ്ഠ കഴിഞ്ഞു ശിവരാത്രിക്കുമുൻപ് അരുവിപ്പുറത്ത് മടങ്ങി എത്തുകയും ചെയ്തു. ഈ യാത്രയിൽ സ്വാമി സഞ്ചരിച്ച പല സ്ഥലങ്ങളിലുമുള്ള പ്രധാന യോഗ്യന്മാരെ യോഗത്തിൽ അംഗങ്ങളായി ചേർക്കുകയും, ചില ദിക്കിൽ മതസംബന്ധമായ ചില സ്ഥാപനങ്ങൾ ആരംഭിക്കാൻ ഏർപ്പാടു ചെയ്കയും ചെയ്തു.
	
	സ്വാമി 1079-ആണ്ടുമുതൽ ആചാരപരിഷ്കരണവിഷയത്തിൽ ദൃഷ്ടിവെക്കയും സ്വജങ്ങളുടെ ഇടയിൽ താലികെട്ട് മുതലായ അനാവശ്യ അടിയന്തിരങ്ങളെ നിർത്തൽ ചെയ്യുവാനും ഒരു പുതിയ വിവാഹരീതിയും ചടങ്ങുകളും ഏർപ്പെടുത്തി പ്രചാരപ്പെടുത്തുവാനും ആരംഭിച്ചു. ഈ സംഗതികൾക്കായി സ്വജനങ്ങലുടെ പല മഹാസഭകളിലും സ്വാമിതന്നെ സന്നിഹിതരായിരുന്നു ജനങ്ങളെ ഗുണദോഷിക്കുകയും ഈ അഭിപ്രായത്തെ എസ്. എൻ. ഡി. പി. യോഗം(ശ്രീ നാരായണ ധർമ്മ പരിപാലനയോഗം) മൂലമായും മറ്റും പ്രകാരത്തിലും പ്രചരിപ്പിക്കുകയും ചെയ്തു. അതിന്റെ ഫലമായി ഈഴവരുടെ ഇടയിൽ മിക്ക സ്ഥലങ്ങളിലും കെട്ടു കല്യാണം വേഗത്തിൽ നിന്നു പോവുകയും ഒരു പുതിയ വിവാഹരീതി നടപ്പാകയും ചെയ്തു.
	
	മദിരാശി ഹൈക്കോർട്ടു ജഡ്ജി സദാശിവയ്യരവർകൾ തിരുവിതാംകൂറിൽ ചിഫ് ജസ്റ്റിസ് ആയിരിക്കുമ്പോൾ ഒരു വിധി കല്പിച്ചതിൽ ഇങ്ങനെ പ്രസ്താവിച്ചിരിക്കുന്നു.
	
	``I hope i might be pardoned for expressing in conclusion my very great satisfaction that through the efforts of the venerable Asan of the Ezhava Community and his Ezhava Samajam, most desirable reforms (or rather the relinquishment of the medieval pernicious customs and conventions which have outlived their original usefulness and which are unsuited to the needs of a progressive community) are taking place among the Ezhavas without the necessity at present to the resort to the legislature. Allude especially to the fast-dying customs of polygamy and polyandry (through restricted to the case of woman being the common wife of brothers) the now unmeaning Thalikettu or Minnukettu ceremoney, the conniving by the Society at the loss of virginity by an unmarried girl remaining in her mother's house and so on. C.A. No: 46 \& 47 of 1083.
	
	ഈഴവ സമുദായത്തിലെ വന്ദ്യനായ ആശാന്റെയും അദ്ദേഹത്തിന്റെ കീഴിലുള്ള ഈഴവ സമാജത്തിന്റെയും ഉദ്യമത്താൽ ഏറ്റവും സ്പൃഹണീയമായ പരിഷ്കാരങ്ങൾ (അതായത് ആദ്യകാലത്തെ ഉപയോഗം നശിച്ചു വർദ്ധമാനമായ സമുദായത്തിന്റെ സ്ഥിതിക്ക് അനുചിതവും ദോഷകരവുമായ വിധത്തിൽ ശേഷിച്ചിരുക്കുന്ന ഇടക്കാലത്തെ ആചാരങ്ങളേയും നടപടികളേയും നിർത്തൽ ചെയ്യുന്ന ഏർപ്പാട്) നിയമ നിർബന്ധം കൂടാതെ തന്നെ ഇപ്പോൾ ഈഴവരുടെ ഇടയിൽ നടന്നുവരുന്നതിനെപ്പറ്റി എനിക്കുള്ള സന്തോഷത്തെ പ്രസ്താവിക്കുന്നത് ക്ഷന്തവ്യമായിരിക്കുമെന്ന് വിശ്വസിക്കുന്നു. ശീഘ്രത്തിൽ നാമാവശേഷമായിക്കൊണ്ടിരിക്കുന്ന ബഹുഭാര്യത്വത്തേയും ബഹുഭർതൃത്വത്തേയും (ബഹുഭർത്തൃത്വത്തിൽ അനേക സഹോദരന്മാർക്കുകൂടി ഒരു ഭാര്യ എന്നുള്ള ഒരു ക്ലിപ്തമുണ്ടെന്നിരുന്നാലും) ഇപ്പോൾ അർത്ഥശൂന്യമായിത്തിർന്നിരിക്കുന്ന താലികെട്ട് അല്ലെങ്കിൽ മിന്നുകെട്ടു കല്യാണത്തേയും മറ്റുമാണു ഞാൻ പ്രത്യേകം ഇവിടെ സൂചിപ്പിക്കുന്നത്. ഈ മിന്നുകെട്ടുകൊണ്ട് വിവാഹം കഴിയാതെ അമ്മയുടെ വീട്ടിൽ താമസിക്കുന്ന പെൺകുട്ടിയുടെ ബ്രഹ്മചര്യഭംഗത്തിൽ സമുദായം കടാക്ഷിക്കയാണ്. (തി.ല.റി. 24-ആം വാള്യം 157 മുതൽ 168 വരെ പേജുകൾ നോക്കുക 1083-ൽ കൃ.അ.നമ്പർ 46-ഉം 47-ഉം).
	
	ഇതിൽനിന്നു സ്വാമിയുടെ ആചാര പരിഷ്കരണ സംബന്ധമായ ഏർപ്പാടുകൾക്കു സിദ്ധിച്ചിട്ടുള്ള വിദ്വൽസമ്മതിയും അഭിനന്ദവും വെളിവാകുന്നതാണ്.
	
	ഈ കാലത്തും സ്വാമി സാഹിത്യ സംബന്ധമായ ശ്രമങ്ങളിൽ നിന്നു വിരമിച്ചിരുന്നില്ല. സ്വാമിയുടെ പരിപക്വമായ ജ്ഞാനാനന്ദഭൂതികളെ സംഗ്രഹിച്ചെഴുതിയിട്ടുള്ള ``ആത്മോപദേശശതകം" എന്ന മണിപ്രവാള പദ്യഗ്രന്ഥം പുറത്തുവന്നത് ഈ അവസരത്തിലാണ്.
	
	\begin{myverse}
		``അയലു തഴപ്പതിനായി പ്രയത്നം \\
		നയമറിയും നരനാചരിച്ചിടേണം.''\\[1\baselineskip]
		``അവരവരാത്മസുഖത്തിനാചരിക്കു-\\
		ന്നവ,യപരന്റെ സുഖത്തിനായ്‌വരേണം.''\\[1\baselineskip]
		``പലമതസാരവുമേകമെന്നു പാരാ\\
		തുലകിലൊരാനയിലന്ധരെന്ന പോലെ\\
		പലവിധയുക്തിപറഞ്ഞു പാമരന്മാ-\\
		രലവതു കണ്ടലയാതിരുന്നിടേണം.''
	\end{myverse}
	
	എന്നിങ്ങനെയുള്ള തത്വരത്നങ്ങൾ ആ മണിപ്രവാളമാലയിൽ ധാരാളമാണ്.
	
	\newpage
	
	\chapter*{അദ്ധ്യായം ഒൻപത്}
	
	ഈ കാലത്ത് സ്വാമി കൂടെക്കൂടെ കുറ്റാലം, പാപനാശം മുതലായ തീർഥസ്ഥലങ്ങളിൽ പോയി വിശ്രമിച്ചിരുന്നു. 1079-ാമാണ്ട് സ്വാമി വർക്കല ഇപ്പോൾ ശിവഗിരി മഠം സ്ഥാപിച്ചിരുക്കുന്ന കുന്നിനു സമീപം ഒരു ദിക്കിൽ പതിവായി ചെന്നിരിക്കുകയും ഒരു കുടിലുകെട്ടി അതിൽ കുറേനാൾ താമസിക്കുകയും ചെയ്തിരുന്നു. അതിനു ചുറ്റും വഴുതിന, പയർ, കത്തിരി, വെണ്ട മുതലായ സസ്യങ്ങൾ കൃഷി ചെയ്യിക്കയും ചെയ്തുകൊണ്ടിരുന്നു. ആ കൃഷിസ്ഥലത്തിന്റെ തെക്കു വശത്തായി ഒരു കുന്നുണ്ടായിരുന്നത് ആരുടേയും പേരിൽ പതിഞ്ഞിട്ടില്ലെന്നു മനസിലാവുകയാൽ തന്റെ ഇരുപ്പു സ്വാമി ക്രമേണ ആ സ്ഥലത്തേക്കു മാറ്റി. കുന്നിന്റെ മുകളിൽ ഒരു പർണ്ണശാല കെട്ടി മിക്കവാറും സ്ഥിരമായി തന്നെ താമസിച്ചു എന്നു പറയാം.
	
	മുൻപ് അരുവിപ്പുറത്ത് എന്ന പോലെ പലസ്ഥലത്തുനിന്നും ജനങ്ങൾ അവിടെ വന്നുകൂടാൻ തുടങ്ങി. കുന്നിന്റെ മുകളിൽ സ്വാമിയുടെ പർണ്ണശാല ഇരുന്ന സ്ഥലത്താണ് ഇപ്പോൾ ശിവപ്രതിഷ്ഠ ചെയ്തിരിക്കുന്നത്. ഇവിടേയും സ്വാമി ക്രമേണ മഠങ്ങളും ക്ഷേത്രങ്ങളും കെട്ടാൻ ആരംഭിച്ചു. കുന്നു തന്റെ പേരിൽ പതിപ്പിക്കുകയും ദാനമായും മറ്റും കിട്ടിയ സമീപത്തുള്ള സ്ഥലങ്ങൾ അതോടു ചേർക്കുകയും സ്ഥലത്തിനു ശിവഗിരി എന്നു പേർകൊടുക്കുകയും ചെയ്തു. ഈ കൊല്ലം ആദ്യമാണ് സിവിൽ കോടതികളിൽ ഹാജരാകേണ്ട നിർബന്ധത്തിൽ നിന്നു സ്വാമിയെ എസ്. എൻ. ഡി. പി യോഗത്തിന്റെ അദ്ധ്യക്ഷന്റേയും സമുദായ ഗുരുവിന്റെയും നിലയിൽ തിരുവിതാംകൂർ ഗവണ്മെന്റിൽ നിന്നും ഒഴിവാക്കിയത്.
	
	1080 ധനുമാസത്തിൽ എസ്. എൻ. ഡി. പി യോഗത്തിന്റെ രണ്ടാമത്തെ വാർഷികയോഗം ഒരു വ്യവസായ പ്രദർശനത്തോടുകൂടി കൊല്ലത്തുവച്ചു നടന്നു. യോഗത്തിന്റെ സ്ഥിരം പ്രസിഡണ്ടിന്റെ നിലയിൽ അലംകരിക്കപ്പെട്ട ഒരു ക്യാബിൻബോട്ടിൽ സ്വാമിയെ ശിവഗിരിയിൽ നിന്നു കൊല്ലത്തേക്കു കൊണ്ടുപോകുവാൻ ജനങ്ങൾ ഉത്സാഹപൂർവ്വമായ ഒരുക്കം കൂട്ടി. പ്രകൃത്യാ ആഡംബരവിമുഖനായ സ്വാമി ആ ബോട്ടിൽ കയറിപ്പോകാതെ തൽക്കാലം ജനങ്ങൾക്ക് വലിയ ആശാഭംഗത്തെയാണ് ഉണ്ടാക്കിയതെങ്കിലും ആ മഹായോഗം കൂടിയിരുന്ന സന്ദർഭത്തിൽ തന്റെ അപ്രതീക്ഷിതമായ സാന്നിദ്ധ്യത്താൽ സഭയെ അലങ്കരിക്കുകയും ജനങ്ങളെ സവിശേഷം സന്തോഷിപ്പിക്കുകയും ചെയ്തു.
	
	\newpage
	
	\chapter*{അദ്ധ്യായം പത്ത്}	
	
	സ്വാമിയുടെ പേർ മുൻപുതന്നെ മലബാറിൽ ശ്രുതിപ്പെട്ടിരുന്നു. അതുകൊണ്ടും ശ്രീനാരായണ ധർമ്മ പരിപാലന യോഗത്തിന്റെ സമുദായ സംബന്ധമായ പ്രവർത്തികൾ നിമിത്തവും അവിടുത്തെ സ്വജനങ്ങൾക്കു സ്വാമിയെ സംബന്ധിച്ച സംഗതികളിൽ പൂർവധികമായ ശ്രദ്ധയും ബഹുമാനവും ഉണ്ടായി. 1081-ാമാണ്ടു തലശ്ശേരിയിൽ സ്വജനങ്ങളുടെ ആരാധനയ്ക്കായി സ്വാമിയെക്കൊണ്ട് ഒരു ക്ഷേത്രം പ്രതിഷ്ഠിപ്പിക്കണമെന്നു ചിലരുടെ ഇടയിൽ ഒരു ആലോചന ഉണ്ടായി. ആ ആണ്ടു കുംഭമാസത്തിൽ അരുവിപ്പുറത്തെ ശിവരാത്രി കഴിഞ്ഞു സ്വാമി തിരുവിതാംകൂറിൽ കോട്ടയം മുതലായ സ്ഥലങ്ങളിൽ സഞ്ചരിക്കയും അവിടെ കുമരകം ഈഴവസമാജം വക ഒരു ക്ഷേത്രത്തിന്റെ പ്രാരംഭകൃത്യം നടത്തുകയും ആ പ്രദേശങ്ങളിൽ ഈഴവരുടെ വകയായി ഉണ്ടായിരുന്ന അനേകം പുരാതന ദുർഗ്ഗാക്ഷേത്രങ്ങളിലെ ജന്തുഹിംസ നിർത്തൽ ചെയ്കയും ചെയ്തു. സ്വാമി ആ സ്ഥലത്ത് ഈ യാത്ര ചെയ്തത് ആദ്യമായിട്ടായിരുന്നു. ഈ യാത്രയിൽ അവിടെ രോഗപീഡിതരായും മറ്റും അനേകം ആളുകൾ വന്നു കൂടിയിരുന്നു. അപ്പോൾ ചില അത്ഭുതസംഭവങ്ങൾ നടന്നതായി അറിയുന്നു. കോട്ടയത്തു നിന്നും സ്വാമി പെരിങ്ങോട്ടുകരയ്ക്കു പോയി. അവിടെ അതിനു മുൻപു തന്നെ ഒരു ക്ഷേത്രവും മഠവും ആരംഭിക്കയും അവയുടെ ഉപയോഗത്തിനായി ഏതാനും സ്വത്തുകൾ സ്വാമിയുടെ പേർക്കു ദാനമായി എഴുതിവാങ്ങുകയും ചെയ്തിരുന്നു. അവിടെ താമസിക്കുമ്പോൾ തലശ്ശേരിയിൽ നിന്നും സ്വജങ്ങളിൽ ചില മാന്യന്മാർ വന്നു സ്വാമിയെ ക്ഷണിച്ചു അങ്ങോട്ടു കൂട്ടിക്കൊണ്ടു പോയി. അവിടത്തെ തിയരിൽ ചിലർ ബ്രഹ്മസമാജത്തിൽ ചേർന്നിരുന്നതുകൊണ്ട് ദീർഘകാലമായി അവരുടെ ഇടയിൽ ഉണ്ടായിരുന്ന കക്ഷിവഴക്കുകൾ സ്വാമി പരഞ്ഞുതീർത്ത് അവരിൽ ഐക്യമത്യം വർദ്ധിപ്പിച്ചു. ഈ ആദ്യത്തെ യാത്രയിൽ തന്നെ സ്വാമി വടക്കെ മലബാറിലുള്ള സ്വജങ്ങളുടെ മുഴുവൻ ഭക്തിക്കും ബഹുമാനത്തിനും സ്നേഹത്തിനും പാത്രമായിത്തീരുകയും തലശ്ശേരി ജഗന്നഥക്ഷേത്രത്തിനു സ്ഥലം നിശ്ചയിച്ചു കുറ്റി തറയ്ക്കുകയും ചെയ്തു. മടക്കത്തിൽ കോഴിക്കോട്ട് ഒന്നു രണ്ടു ദിവസം താമസിക്കയും പെരിങ്ങോട്ടുകര, പറവൂർ ഈ സ്ഥലങ്ങളിൽ കൂടി യാത്ര ചെയ്തു മീനമാസ മദ്ധ്യത്തോടു കൂടി ശിവഗിരിയിൽ വന്നു ചേരുകയും ചെയ്തു. ഈ യാത്രയിൽ പറവൂർ മുതലായ സ്ഥലങ്ങളിൽ പഴയ താലികെട്ടു നിർത്തൽ ചെയ്യുവാനും പുതിയ വിവാഹരീതി പ്രചാരപ്പെടുത്താനും സ്വാമി ഏർപ്പാടു ചെയ്തു.
	
	\newpage
	\chapter*{അദ്ധ്യായം പതിനൊന്ന്}
	
	പിന്നെ രണ്ടു കൊല്ലത്തക്ക് (83-ാമാണ്ടുവരെ) സ്വാമിമിക്കവാറും ശിവഗിരിയിൽ തന്നെ വിശ്രമിക്കയായിരുന്നു. ഈ കാലത്താണ് സ്വാമി അവിടത്തെ പഴയമഠത്തിന്റെ പണിപൂർത്തിയാക്കുകയും ഒരു സംസ്കൃതപാഠശാല ഏർപ്പെടുത്തുകയും മറ്റും ചെയ്തത്. ചീഫ് ജസ്റ്റീസ് സദാശിവയ്യർ മുതലായ പല മഹത്തുക്കളും ഈ അവസരത്തിൽ സ്വാമിയെ ചെന്നു സന്ദർശിച്ചു തൃപ്തിപ്പെട്ടിട്ടുണ്ട്. സ്വജനങ്ങളിൽ വടക്കും തെക്കുമുള്ള പലപ്രധാന യോഗ്യന്മാരും ഈ അവസരത്തിൽ കൂടെക്കൂടെ അവിടെ വന്നു സ്വാമിയെ സന്ദർശിച്ചുകൊണ്ടിരുന്നു. ക്രമേണ ശിവഗിരി മഠത്തിൽ ഒരു പൊതുക്ഷേത്രവും ഒരു ഉയർന്നതരം സംസ്കൃതവിദ്യാമന്ദിരവും സ്ഥാപിക്കണമെന്നും ശിവഗിരി മഠത്തെ തന്റെ മതസംബന്ധമായ സ്ഥാപനങ്ങളുടെ തലസ്ഥാനമാക്കണമെന്നും ഉള്ള അഭിപ്രായത്തെ സ്വാമി വെളിപ്പെടുത്തി. 1083 ചിങ്ങം 13-ാം തിയതി സ്വാമി പരസ്യമായി ഇതിലേക്ക് സ്വജനങ്ങളിൽ നിന്നും ധനസഹായം അപേക്ഷിച്ച് ഒരു ചെറിയ വിജ്ഞാപനം പ്രസിദ്ധപ്പെടുത്തി.
	
	ഈ സമയത്ത് തലശ്ശേരിയിൽ സ്വജനങ്ങളുടെ വക ക്ഷേത്രത്തിന്റെ പണികൾ പൂർത്തിയായി കൊണ്ടിരുന്നു. കോഴിക്കോട്ടും ഒരു ക്ഷേത്രം സ്ഥാപിക്കുന്നതിനുള്ള ആലോചനകൾ എല്ലാം കഴിച്ച് അതിന്റെ പ്രാരംഭക്രിയകൾക്ക് ഒരുക്കം കൂട്ടിക്കൊണ്ടിരുന്നു. അതു സംബന്ധിച്ചു സ്വാമിയെ അവിടത്തേക്കു ക്ഷണിപ്പാൻ മ. രാ. രാ. കല്ലിങ്കൽ രാരിച്ചൻ മൂപ്പൻ മുതലായ മാന്യന്മാർ ശിവഗിരിയിൽ വന്നിരുന്നു. അവരെ എല്ലാം മുൻകൂട്ടി മടക്കി അയച്ചിട്ട്, സ്വാമി 1088 തുലാം 24-ാം തിയതി രാത്രി ഏതാനും അനുയായികളുമായി യാത്രപുറപ്പെട്ടു. കൂടെയുള്ളവർ നിർബന്ധിക്കയാലാണ് താൻ പുറപ്പെടുന്നതെന്നും ഉള്ളിൽ അശേഷം ഉത്സാഹം തോന്നുന്നില്ലെന്നും മറ്റും പറഞ്ഞു യാത്രാരംഭത്തിൽ സ്വാമി വളരെ മടിച്ചു. 27-ാം തിയതി വൈകുന്നേരം ആലുവായിൽ എത്തി. പിറ്റേദിവസത്തെ മെയിൽ വണ്ടിയിൽ കോഴിക്കോട്ടേക്കു പോവാൻ വിചാരിച്ചുകൊണ്ട് പരിജന സഹിതം രാത്രി തീവണ്ടി സ്റ്റേഷനുസമീപം ഒരു കെട്ടിടത്തിൽ താമസിച്ചു. ആ കെട്ടിടം ഇപ്പോൾ ആലുവാ അദ്വൈതാശ്രമം വക ഒരു മഠമാക്കിയിരിക്കയാണ്. രാത്രി 12 മണിക്കു സ്വാമിക്കു അവിടെ വച്ച് അതിസാരത്തിന്റെ ലക്ഷണം ആരംഭിച്ചു. ഒരു മണിക്കൂർ കഴിഞ്ഞു ഛർദ്ദികൂടി തുടങ്ങുകയും ഭയങ്കരമായ വിഷൂചികയാണെന്നു വെളിപ്പെടുത്തുകയും ചെയ്തു. ആരംഭം മുതൽ തന്നെ സ്ഥലത്തെ അപ്പാത്തിക്കിരിയെക്കൊണ്ട് ഇംഗ്ലീഷ് ചികിത്സ വളരെ ജാഗ്രതയോടെ ചെയ്യിക്കയും സകലവിധമായ ശുശ്രൂഷകളും യാതൊരു ന്യൂനതയും കൂടാതെ കൂടെയുള്ളവർ നടത്തുകയും ചെയ്തിരുന്നു. എന്നാൽ ദീനം ഭയങ്കരമാംവിധം വർദ്ധിച്ചു കൊണ്ടുതന്നെയിരുന്നു. പിറ്റേദിവസം പകൽ 12 മണിയോടു കൂടി പ്രജ്ഞ അശേഷം കെടുകയും നാഡി നിന്നുപോകയും ശരീരം തണുത്തുമരവിച്ച്പോകുകയും ചെയ്തു. ഇതിനിടയിൽ കൂടെയുണ്ടായിരുന്ന ശിഷ്യന്മാരിൽ `കൊച്ചുമായിറ്റി' ആശാൻ എന്ന ഒരാൾക്കുകൂടി ദീനം ആരംഭിച്ച് അദ്ദേഹം മരിച്ചു കഴിഞ്ഞിരുന്നു. ഡോക്ടറുടെ അഭിപ്രായം അദ്ദേഹത്തിന്റെ സുഖക്കേടു ശമിക്കുമെന്നും സ്വാമിയുടെത് അസാധ്യം എന്നും ആയിരുന്നു. ഈ സമയത്ത് കോഴിക്കോട്ടുതീവണ്ടിയാഫീസിലും സമീപപ്രദേശങ്ങളിലും മെയിൽ വണ്ടിയിൽ സ്വാമി വന്നിറങ്ങുന്നതു കാണ്മാനും എതിരേൽക്കാനുമായി കൂടിയിരുന്ന പുരുഷാരത്തിനു കണക്കില്ലായിരുന്നു. ആലുവായിൽ സ്വാമിയുടെ സകലശുശ്രൂഷകളും അവസാനിപ്പിച്ചു കണ്ണീർവാർത്തുകൊണ്ട് വിഷണ്ണന്മാരായി നിൽക്കുന്ന ശിഷ്യന്മാരുടെയും പരിജനങ്ങളുടേയും ഈ സമയത്തെ ഹൃദയസ്ഥിതി പറഞ്ഞറിയിപ്പാൻ കഴിയാത്തതായിരുന്നു.
	
	\newpage
	\chapter*{അദ്ധ്യായം പന്ത്രണ്ട്}
	
	ഏതാനും ദിവസങ്ങൾകൊണ്ടു സ്വാമി ആ മഹാരോഗത്തിൽ നിന്നും, അതിനെ തുടർന്നുണ്ടായ ഉപദ്രവങ്ങളിൽ നിന്നുമെല്ലാം അദ്ഭുതകരമാംവണ്ണം രക്ഷപ്പെട്ടു. മുമ്പൊരിക്കൽ യോഗ ശീലിച്ചുകൊണ്ടിരുന്ന കാലത്തു നെയ്യാറ്റുംകരയ്ക്കടുത്ത അരുമാനൂർ എന്ന ഗ്രാമത്തിൽ വച്ച് സ്വാമിക്ക് ഇതുപോലെ വിഷൂചിക ബാധിച്ചു. അന്നു സകല ചികിത്സകളും നിറുത്തി സ്വാമി സ്വസ്ഥമായിരുന്നപ്പോൾ രോഗം പെട്ടെന്നുമാറി സുഖം വന്നതായി കേട്ടിട്ടുണ്ട്. പ്രസിദ്ധ ഡാക്ടറായ റാവു സാഹിബ് കെ. കൃഷ്ണൻ അവർകൾ സ്വാമിയെ ആലുവായിൽ നിന്ന് ചികിത്സക്കായി പാലക്കാട്ടുകൊണ്ടുപോയി താമസിപ്പിച്ചിരുന്ന അവസരത്തിൽ അവിടത്തെ സ്വജനങ്ങളുടെ വകയായി സ്ഥലത്ത് ഒരു ക്ഷേത്രവും സഭയും ഏർപ്പെടുത്തുന്നതിനെപ്പറ്റി ചില ആലോചനകൾ എല്ലാം നടന്നു. ധനു 8-ാം തിയതി സ്വാമി അവിടെ ഒരു ക്ഷേത്രത്തിനായി ആഘോഷപൂർവ്വം കുറ്റി തറക്കുകയും ചെയ്തിട്ടുണ്ട്. അടുത്ത ദിവസം സ്വാമി കോഴിക്കോട്ടെത്തി. സ്വാമിയെ സ്വീകരിപ്പാൻ കോഴിക്കോട്ടെ സ്വജനങ്ങൾ ചെയ്തിരുന്ന ഏർപ്പാടുകൾ സ്തുത്യർഹങ്ങളായിരുന്നു. 12-ാ൦ തിയതി കാലത്ത് അതികേമമായ ആഘോഷത്തോടുകൂടി സ്വാമി കോഴിക്കോട്ടു ശ്രീകണ്ഠേശ്വരം ക്ഷേത്രത്തിനു കുറ്റിതറച്ചു. 15-ാ൦ തിയതി രാവിലെ ആനി ഹാളിൽ (Anne Hall) വച്ച് സ്ഥലം തിയോസഫി സഭയിലെ അംഗങ്ങളും വിദ്വാന്മാരുമായി ഏതാനും ബ്രാഹ്മണരും നായന്മാരും ചേർന്നു സ്വാമിക്ക് ഒരു മംഗളപത്രം സമർപ്പിച്ചു. അവർക്കു സ്വാമിയെപ്പറ്റി തോന്നിയ ഭക്തി ബഹുമാനങ്ങൾ ടി മംഗളപത്രത്തിൽ നിന്ന് താഴെ പകർത്തുന്ന വാക്യങ്ങളിൽ നിന്നു വെളിവാകുന്നതാണ്.
	
	\begin{myverse}
		ഏവം സദ്ധർമ്മകർമ്മാചരണ വിഷയിണിം\\
		\hspace*{2cm}ബുദ്ധി മസ്മാകമത്രാ\\
		ധാതും കാരുണ്യപൂർവ്വം നിജനിലയമപാ-\\
		\hspace*{2cm}ഹായ ചാഭ്യാഗമദ്യത്\\
		ശ്രീമന്നാരായണാഖ്യഃ ശിവഗിരിനിലയാ-\\
		\hspace*{2cm}ധിശ്വരോ ദേശികോയം\\
		തസ്മത് സന്തുഷ്ടാചിത്താ വയമപി\\
		\hspace*{2cm}തനുമഃ സ്വാഗതംഃ വന്ദനം ച.
	\end{myverse}
	
	We.....recognising in you a born leader of men, a genuine descendant of the ancient saints of our mother land, a true Bramhana soul sent out by the guardians of humanity for the uplifting and redemption of a community whose spiritual interest those who call themselves highcaste have grown so sadly oblivious.\\
	(സ്വാമി മനുഷ്യവർഗ്ഗത്തിന്റെ - ബ്രാഹ്മണാത്മാവാകുന്നു.)
	
	ഈ യാത്രയിൽ സ്വാമി തലശ്ശേരിവരെ പോകയും അവിടെ ക്ഷേത്രസംബന്ധമായും ``ജ്ഞാനോദയയോഗം" സംബന്ധിച്ചും തൽപ്രവർത്തകന്മാരുടെയും പൊതുജങ്ങളുടെയും ഇടയിൽ നടന്നിരുന്ന ചില തർക്കങ്ങളേയും കുഴപ്പങ്ങളേയും പറഞ്ഞുതീർക്കുകയും ചെയ്തു. സ്വാമി സ്വല്പദിവസംകൂടി വടക്കേമലയാളത്തിലുള്ള പല ചെറിയ പട്ടണങ്ങളും നഗരങ്ങളും സന്ദർശിച്ചുകൊണ്ടു താമസിച്ചു. മകരമാസത്തിൽ സ്വാമി കണ്ണൂർ സ്വജങ്ങളുടെ വകയായ ഒരു ക്ഷേത്രത്തിന് കുറ്റി തറക്കുകയും അവിടെനിന്നു മംഗലാപുരത്തു വില്ലവർ എന്ന തുളുതിയരുടെ വകയായി ഒരു ക്ഷേത്രം സ്ഥാപിക്കുന്നതിന്റെ പ്രാരംഭപ്രവൃത്തികൾക്കായി അവിടം വരെ പോവുകയും ചെയ്തു. 1083 കുംഭം 1-ാം തിയതി തലശ്ശേരി ജഗന്നാഥക്ഷേത്രത്തിലെ പ്രതിഷ്ഠാകർമ്മം നടത്തുകയും ക്ഷേത്രത്തിലെ തന്ത്രാവകാശം ശാശ്വതമായി ശിവഗിരിമഠത്തിൽ ഇരിക്കത്തക്കവണ്ണം ജ്ഞാനോദയ യോഗക്കാരിൽനിന്ന് ഉടമ്പടി എഴുതിവാങ്ങി ഏർപ്പാടുചെയ്യുകയും ചെയ്തു. മടക്കത്തിൽ സ്വാമി കൊച്ചിയിൽ, കുളമ്പടി എന്ന സ്ഥലത്ത് ഒരു പഴയ ദുർഗ്ഗാക്ഷേത്രം പുതുക്കി പ്രതിഷ്ഠിക്കയും അവിടെ നടന്നുവന്ന പൂരംതുള്ളൽ, കുരുതി മുതലായ നിഷിദ്ധാചാരങ്ങളെ നിർത്തൽചെയ്യുകയും ചെയ്തു. ആ ആണ്ടു മേടം, ഇടവം ഈ മാസങ്ങളിൽ ശിവഗിരിയിലെ ധർമ്മകാര്യങ്ങൾ സംബന്ധിച്ചു സ്വാമി കൊല്ലം മുതലായ താലൂക്കുകളിൽ സഞ്ചരിച്ചു.
	
	\newpage
	\chapter*{അദ്ധ്യായം പതിമൂന്ന്}
	
	1084 ചിങ്ങം 26-ാം തീയതി (ചതയം) സ്വാമിയുടെ 53-ാമത്തെ ജന്മനക്ഷത്രാഘോഷമായിരുന്നു. അന്നു ശിവഗിരിയിൽ സ്വജനങ്ങളുടെ ഒരു യോഗം കൂടുകയും സ്വാമി ``ശാരദാമഠ"ത്തിന് അടിസ്ഥാനക്കല്ലുവയ്ക്കുകയും ചെയ്തു. ``ജനനീനവരത്നമഞ്ജരി" എന്ന മധുരഗംഭീരമായ ചെറിയ സ്തോത്രകൃതി സ്വാമി ഇക്കാലത്തു ശിവഗിരിയിൽ വെച്ചു നിർമ്മിച്ചതാണ്. ശിവഗിരി മഠത്തിൽ അപ്പോൾ മഠം വക പണികൾക്കു പുറമേ മതസംബന്ധമായ ശാസ്ത്രപഠനങ്ങൾ, പാരായണം മുതലായവും നിയമേന നടന്നുകൊണ്ടിരുന്നു. പാണ്ഡിത്യമുള്ള പല ശിഷ്യന്മാരും ഗീത, ഉപനിഷത്തുക്കൾ,വസിഷ്ഠം സൂതസംഹിത മുതലായ പ്രൗഢഗ്രന്ഥങ്ങളുടെ അർത്ഥങ്ങൾ സ്വാമി മുഖേന കേട്ടു ഗ്രഹിക്കുകയും അവയെ ചർച്ച ചെയ്കയും ചെയ്തുകൊണ്ടിരുന്നു. ക്ഷേത്രസംബന്ധമായ തന്ത്രങ്ങൾ സമുദായികമായ സംസ്കാരകർമ്മങ്ങൾ മുതലായവ അഭ്യസിപ്പിക്കുന്നതിനു മഠത്തിൽ ഏർപ്പാടുകൾ ചെയ്തിരുന്നു. സ്വജനങ്ങളുടെ വിവാഹം , അപരകർമ്മങ്ങൾ, ശ്രാദ്ധം, ക്ഷേത്രപിണ്ഡം മുതലായവ മഠത്തിൽനിന്നും വൈദികന്മാരെ അയച്ചു ശാസ്ത്രരീത്യാ ചെയ്തുകൊടുത്തിരുന്നു. ഇടവമാസത്തിൽ എറണാകുളത്തുവച്ചു കൂടിയ എസ്.എൻ.ഡി.പി യോഗത്തിന്റെ 6-ാമതു വാർഷിക പൊതുയോഗത്തിൽ മതസംബന്ധമായും ആചാരപരിഷ്കരണസംബന്ധമായുമുള്ള പ്രവൃത്തികളിൽ യോഗം ദൃഷ്ടിവെക്കേണ്ട സാരമായ ചില സംഗതികളെപ്പറ്റി സ്വാമി എഴുതി അറിയിക്കയും ആ വിഷയങ്ങളിൽ പൊതുജനങ്ങളെ അതുമൂലം ഊർജ്ജിതപ്പെടുത്തുകയും ചെയ്തു. ഇതിന് ഏതാനും ദിവസങ്ങൾക്കു മുമ്പ് വടക്കേമലയാളത്തെ സ്വജനങ്ങളിൽ അനേകായിരം ആളുകൾ ആണ്ടുതോറും കൊട്ടിയൂർ എന്നുകൂടി പറയാറുള്ള തൃച്ചംബരത്തു ക്ഷേത്രത്തിൽ ഇളനീരഭിഷേകത്തിനുപോയി അനവധി പണം വൃഥാ ചെലവുചെയ്കയും കാണിക്ക ഇടുകയും ചെയ്യുന്ന പതിവു നിറുത്തി അവരുടെ ഭക്തിവിശ്വാസങ്ങളെ തലശ്ശേരി ജഗന്നാഥ ക്ഷേത്രത്തിലേക്കു തിരിച്ചുവിടത്തക്കവണ്ണം വേണ്ടതാലോചിക്കുകയും തന്റെ ആ ആഗ്രഹം ജനങ്ങൾ അറിവാനായി പത്രദ്വാരാ പ്രസിദ്ധപ്പെടുത്തുകയും ചെയ്തു. അതിന്റെ ഫലമായി ഇടവം 29-ാം തീയതി ജഗന്നാഥക്ഷേത്രത്തിൽ അതികേമമായ വിധത്തിൽ അഭിഷേകോത്സവം നടന്നു. ഒരു വലിയ സംഖ്യ നടവരവും ഉണ്ടായി. അക്കുറി കൊട്ടിയൂർ പതിവുപോലെ പോയ തീയരിൽ പലരും ഭാഗ്യദോഷത്താൽ പുഴയിലെ വെള്ളപ്പൊക്കത്തിൽപ്പെട്ടു മരിച്ചുപോയ വസ്തുത ശേഷമുള്ളവർക്കു സ്വാമിയുടെ വാക്കിൽ വിശ്വാസത്തെ വർദ്ധിപ്പിച്ചു എന്ന സംഗതിയും പ്രസ്താവയോഗ്യമാണ്. സന്യാസികളായ ശിഷ്യൻമാരിൽ യോഗ്യതയുള്ള ചിലരെ സ്വാമി ഈ അവസരത്തിൽ മതം, സദാചാരം മുതലായ വിഷയങ്ങളിൽ പ്രസംഗങ്ങൾ ചെയ്തു ജനങ്ങൾക്ക് അറിവും സന്മാർഗ്ഗനിഷ്ഠയും വർദ്ധിപ്പിപ്പാൻ ചില പ്രത്യേക നിർദ്ദേശങ്ങളോടുകൂടി കേരളത്തിന്റെ പല ഭാഗങ്ങളിലും അയച്ചിരുന്നു. ഈ കൊല്ലത്തിൽ ശിവഗിരിമഠത്തിനു സമീപമായി സ്വാമി മഠത്തിലേക്കു കുറെ ഭൂസ്വത്തുക്കൾ വിലയായും ഒറ്റിയായും വാങ്ങുകയും ചെയ്തു.
	
	1085 ധനു 16-ാം തീയതി ശിവഗിരി വിട്ടു സ്വാമി വീണ്ടും മംഗലപുരത്തേക്ക് തിരുനെൽവേലിവഴിയായി യാത്ര പുറപ്പെട്ട് മദ്ധ്യേ മധുരയിലും മറ്റും ഏതാനും ദിവസം താമസിക്കയും ചെയ്തു. അതിനുശേഷം കോഴിക്കോട്, മയ്യഴി, തലശ്ശേരി, കണ്ണൂർ മുതലായ സ്ഥലങ്ങളിൽ അല്പാല്പം വിശ്രമിച്ചുകൊണ്ടു മംഗലപുരത്തെത്തി. മകരം 3-ാം തിയതി അിടത്തെ (ഇപ്പോൾ തൃപ്പതീശ്വരം എന്നു പറയുന്ന) ക്ഷേത്രത്തിന് ആഘോഷപൂർവ്വം കുറ്റി തറക്കുകയും, മടങ്ങി കുംഭം 2-ാം തീയതിയോടുകൂടി ശിവഗിരിയിൽ എത്തുകയും ചെയ്തു. മേടത്തിൽ വീണ്ടും സ്വാമി മലബാറിലേക്കു പോയി. ആ മാസം 29-ാം തീയതി കോഴിക്കോട്ടു ``ശ്രീകണ്ഠേശ്വരം" എന്ന, സ്വാമി ഒടുവിൽ നാമകരണം ചെയ്ത ക്ഷേത്രത്തിന്റെ കുംഭാഭിഷേകം നടത്തി. മംഗലപുരത്ത് കുറ്റിതറച്ചതായി മേൽ പ്രസ്താവിച്ച ക്ഷേത്രത്തിന്റെ ഷഡാധാരപ്രതിഷ്ഠ സ്വാമി 18-ാം തീയതി നടത്തി. മിഥുനമാസത്തിൽ സ്വാമി നെയ്യാറ്റിൻകര, തിരുവനന്തപുരം ഈ ദിക്കുകളിൽ സഞ്ചരിക്കുകയും, മഠത്തിലേക്ക് ധനാർജ്ജനം ചെയ്യുന്ന വിഷയത്തിൽ യത്നിക്കുകയും ചില വിലപിടിച്ച ദാനാധാരങ്ങൾ എഴുതിവാങ്ങുകയും ചെയ്തു.
	
	1086-തുലാമാസം മുതൽ തിരുവനന്തപുരം നെയ്യാറ്റിൻകര ഈ താലൂക്കുകളിൽ സ്വജനങ്ങളുടെ ആചാര പരിഷ്ക്കരണസംബന്ധമായും മറ്റുമുള്ള ചില സംഗതികൾ സംബന്ധിച്ച് മകരം 5-ാം തീയതി കരുംകുളം എന്ന സ്ഥലത്ത് ഒരു മാന്യ ഈഴവകുടുംബത്തിൽ അതികേമമായി നടന്നുകൊണ്ടിരുന്ന ഒരു താലികെട്ടടിയന്തിരം പന്തലിൽ സ്വാമിചെന്നു പെട്ടെന്ന് കയറി അപ്പോൾതന്നെ മുടക്കുകയും, മേലാൽ ആ സ്ഥലങ്ങളിൽ ആരും താലികെട്ടടിയന്തിരം നടത്തരുതെന്നുപദേശിക്കയും ചെയ്തു.
	
	ശിവഗിരിമഠത്തിൽ സംസ്കൃതപാഠശാല ഈ അവസരത്തിൽ പൂർവ്വാധികം പുഷ്ടിപ്പെടുത്തുകയും സംസ്കൃതഭാഷയ്ക്കു പുറമേ കണക്ക്, ഇംഗ്ലീഷ് മുതലായ പാഠങ്ങൾകൂടി ഏർപ്പെടുത്തുകയും ചെയ്തു. വിദ്യാർത്ഥികളുടെ സംഖ്യ അൻപതിലധികപ്പെട്ടിരുന്നു. ശിവഗിരിയിലെ ശാരദാപ്രതിഷ്ഠ ഈയാണ്ടു മേടത്തിൽ നടത്തണമെന്നു സ്വാമി ആദ്യം ആലോചിച്ചു. ചില ഒരുക്കങ്ങൾ പൂർത്തിയാകാത്തതിനാൽ അപ്പോൾ നടത്താൻ സാധിച്ചില്ല. പിന്നെ 1087 ചിങ്ങത്തിൽ സ്വാമിയുടെ ജന്മനക്ഷത്രം സംബന്ധിച്ചു നടത്തുന്നതു നന്നായിരിക്കുമെന്നു ഭാരവാഹികളിൽ ചിലരുടെ ഇടയിൽ ആലോചന നടന്നു. അന്നും സാധിച്ചില്ല. 1086 ഇടവത്തിൽ സ്വാമി കരുനാഗപ്പള്ളിയിൽ ചില സ്ഥലങ്ങളിൽ സഞ്ചരിച്ചു. ആ അവസരത്തിൽ അവിടെയുള്ള രണ്ടു പ്രധാന ഈഴവകുടുംബക്കാർ തമ്മിൽ നടന്നുവന്ന ഒരു ആപൽകരമായ മത്സരത്തെ രാജിപ്പെടുത്തിയിരുന്നു. 1088 കർക്കടമാസത്തിൽ സ്വാമി ഏതാനും ദിവസം കുറ്റാലത്തുപോയി വിശ്രമിച്ചിരുന്നു. 87 കന്നി 21-ാം തിയതി ശിവഗിരിയിൽ വച്ചു നടന്ന എസ്. എൻ. ഡി. പി. യോഗത്തിന്റെ 8-ാമതു വാർഷികപൊതുയോഗത്തിൽ സ്വാമി സംബന്ധിക്കുകയും ശാരദാപ്രതിഷ്ഠ നടത്തുന്നതിനായി യോഗം മുഖേന ഒരു കമ്മറ്റിയെ നിശ്ചയിച്ചു ചുമതലപ്പെടുത്തുകയും ചെയ്തു. തുലാമാസത്തിൽ സ്വാമി പാപനാശം എന്ന തീർത്ഥസ്ഥലത്തുപോയി കുറേദിവസം വിശ്രമിച്ചു. ഈ യാത്രയിൽ തിരുനൽവേലി, മധുര, തഞ്ചാവൂർ ഈ ഡിസ്ത്രിക്ടുകളിലുള്ള പല പുണ്യസ്ഥലങ്ങളിലും ഏതാനും ദിവസങ്ങൾ താമസിക്കയും അവിടെനിന്നു ധനുമാസത്തിൽ മലബാരിലേക്കു പോവുകയും ചെയ്തു. ജനങ്ങൾ വേണ്ടത്ര ഉത്സാഹിക്കാത്തതിനാൽ ശാരദാപ്രതിഷ്ഠ നടത്താൻ താമസം നേരിടുന്നതിൽ തനിക്കുള്ള വൈമനസ്യംകൊണ്ട് സ്വാമി ശിവഗിരിവിട്ട് സഞ്ചരിക്കയാണെന്ന് ഈ യാത്രയിൽ ഒരു ജനശ്രുതി പരന്നു. സ്വാമിയെ ദീർഘസഞ്ചാരത്തിന് പ്രേരിപ്പിച്ച യഥാർത്ഥമായ ഉദ്ദേശ്യം എന്താണെന്നു അധികംപേർ അറിഞ്ഞിട്ടില്ല. ഒടുവിൽ യോഗം ജനറൽ സെക്രട്ടറി സ്വാമിയെ തലശ്ശേരിയിൽചെന്നു സന്ദർശിച്ചു. ശരദാ പ്രതിഷ്ഠാസംബന്ധമായ അവിടത്തെ ആഗ്രഹങ്ങൾ അറിഞ്ഞും ആജ്ഞാപനങ്ങൾ വാങ്ങിയും മടങ്ങുകയും പ്രതിഷ്ഠ 88 മേടത്തിൽ നടത്താൻ യോഗത്തിന്റെ പ്രധാന ഭാരവാഹികളായ ജനങ്ങളുമായി ആലോചിച്ചും പ്രതിഷ്ഠാക്കമ്മറ്റി മീറ്റിംഗ് കൂടിയും നിശ്ചയിക്കുകയും ചെയ്തു.
	
	\newpage
	\chapter*{അദ്ധ്യായം പതിനാല്}
	
	1088 മകരമാസത്തിൽ സ്വാമി ആലുവായിലും അടുത്ത പ്രദേശങ്ങളിലും സഞ്ചരിച്ചിരുന്നു. മകരം 9-ാം തീയതി കൊച്ചി സംസ്ഥാനത്തു മുനമ്പിനുസമീപം ചെറായി എന്ന സ്ഥലത്ത് സ്വജനങ്ങളുടെ വകയായി ആരംഭിച്ച ക്ഷേത്രത്തിന്റെ പണികൾ കാൺമാനും മറ്റുമായി അവിടെ പോവുകയും സ്ഥലം ``വിദ്യാപോഷിണിസഭ" വകയായി ഒരു മംഗളപത്രം സ്വീകരിക്കയും മതം, സദാചാരം, വിദ്യാഭ്യാസം, വ്യവസായം മുതലായ സകലമാർഗ്ഗങ്ങളിലും സ്വജനങ്ങൾ അഭിവൃദ്ധി സമ്പാദിക്കേണ്ട ആവിശ്യകതയെസൂചിപ്പിച്ച് ഒരു സാരമായ മറുപടി നൽകുകയും ചെയ്തു മടങ്ങി. ക്ഷണമനുസരിച്ചു താമസിയാതെ തന്നെ ടി സ്ഥലത്തേക്കു വീണ്ടും സ്വാമി പോവുകയും മേൽപറഞ്ഞ ക്ഷേത്രത്തിന്റെ പ്രതിഷ്ഠ നടത്തുകയും ``ഗൗരീശ്വരി" എന്നു നാമകരണം ചെയ്കയും ചെയ്തു. ഈ കാലങ്ങളിൽ സ്വാമി ആലുവായിൽ കൂടെക്കൂടെ താമസിക്കുകയും അവിടെ സ്ഥിരമായി എന്തെങ്കിലും ഒരു സ്ഥാപനം ഉണ്ടാക്കാനുള്ള അഭിലാഷത്തെ ഒരു വിധം തന്റെ പ്രവർത്തികളിൽ സൂചിപ്പിക്കയും ചെയ്തു. ഏകദേശം പത്തുകൊല്ലങ്ങൾക്ക് മുൻപ് അരുവിപ്പുറത്തുനിന്നും സ്വാമി വർക്കലെ സഞ്ചരിച്ചിരുന്നതുപോലെ ഈകാലത്ത് ശിവഗിരിയിൽനിന്നും ആലുവായിൽ പൊയ്ക്കൊണ്ടിരുന്നു എന്നു പറഞ്ഞാൽ ഏതാണ്ട് ശരിയായിരിക്കും. ആലുവായിൽനിന്നും സ്വാമി മംഗലപുരത്തേക്ക് ക്ഷണം അനുസരിച്ചു പോകയും കുംഭം 10-ാം തിയതി അവിടത്തെ തീയരുടെ ക്ഷേത്രത്തിന്റെ കുംഭാഭിഷേകം നടത്തുകയും ആ ക്ഷേത്രത്തിന് ``തൃപ്പതീശ്വരം" എന്ന് പേരിടുകയും ചെയ്തു. മംഗലപുരത്തെ തീയർ തൃപ്പതി വിഷ്ണുക്ഷേത്രത്തിൽ ആണ്ടുതോറും വലിയ തീർത്ഥയാത്ര ചെയ്യുകയും അനവധി പണം അവിടെകൊണ്ടുപോയി കാണിക്ക ഇടുകയും ചെയ്യുക പതിവായിരുന്നു. ആ കാണിക്കകൾ തങ്ങളുടെ ക്ഷേത്രത്തിൽ സമർപ്പിച്ചാൽ മതിയെന്നും ദേവന്റെ സാന്നിദ്ധ്യം തൃപ്പതിയിലെപ്പോലെതന്നെ ഇന്ന് ഇവിടെയും ഉണ്ടെന്നും സ്വാമി അവരോടുപദേശിച്ചിരുന്നു. ``തൃപ്പതീശ്വരം" എന്നു ക്ഷേത്രത്തിനു നാമകരണം ചെയ്തതുതന്നെ ആ ഉദ്ദേശത്തിലായിരിക്കണം. തലശ്ശേരി ക്ഷേത്രത്തിനു സ്വാമി ജഗന്നാഥം എന്നുപേർ വിളിച്ചതും ഇതുപോലെ വേറൊരുദ്ദേശത്തോടുകൂടി ആയിരുന്നു. ഇന്ത്യയിലെ പ്രധാന ഹിന്ദുക്ഷേത്രങ്ങളിൽവച്ച് ഒറിസ്സയിലെ പുരി എന്ന സ്ഥലത്തെ ജഗന്നാഥ ക്ഷേത്രം ഏറ്റവും പുരാതനവും മുഖ്യവുമായ ഒന്നാണ്. അവിടെ ജാതിവിചാരം അശേഷം ഇല്ലെന്നുള്ള ഒരു വിശേഷം കൂടിയുണ്ട്. തന്റെ കീഴിലുള്ള ക്ഷേത്രങ്ങൾ എല്ലാം അങ്ങിനെയായാൽകൊള്ളാമെന്ന ആഗ്രഹം സ്വാമിക്കുണ്ട്. അതിനെ സ്വാമി പല അവസരങ്ങളിലും പ്രവർത്തിമൂലം പ്രദർശിപ്പിച്ചിട്ടുമുണ്ട്. തലശ്ശേരി ക്ഷേത്രത്തിൽ ചില വർഗ്ഗക്കാർക്ക് പ്രവേശം അനുവദിക്കുന്നതിന് തീയരിൽ പൂർവാചാര പ്രിയരായ ചിലർക്കു വിരോധം ഉള്ളതായി സ്വാമി അറിഞ്ഞിട്ടുണ്ടായിരുന്നു. ക്ഷേത്രത്തിനു ജഗന്നാഥം എന്നു പേർ കൊടുത്തിട്ടുള്ള വിഷയത്തിൽ തനിക്കുള്ള വിശാലമായ അഭിപ്രായം സ്വാമി ജനങ്ങളെ അറിയിച്ചതാവുന്നു.
	
	മംഗലപുരത്തെ ക്ഷേത്രപ്രതിഷ്ഠയോടുകൂടി സ്വാമിയുടെ ക്ഷേത്രസ്ഥാപനയത്നം കേരളത്തിന്റെ ഒരറ്റത്തുനിന്നും മറ്റേ അറ്റം വരെ എത്തി വേരൂന്നിക്കഴിഞ്ഞു. അതിനുശേഷം അതിന്റെ യഥാക്രമമായ വളർച്ചയേയും നിലനിൽപിനേയും പറ്റിയാണ് സ്വാമിക്ക് ചിന്തിക്കേണ്ടിയുള്ളത്. മുമ്പേതന്നെ തന്റെ മതസംബന്ധമായ സ്ഥാപനങ്ങളുടെ തലസ്ഥാനമാക്കാനുദ്ദേശിച്ചിരിക്കുന്ന ശിവഗിരിയിലെ പ്രതിഷ്ഠയും അതുസംബന്ധിച്ച ഏർപ്പാടുകളും എല്ലാം അതിന്റെ ഗൗരവത്തിനു തക്കവണ്ണം കേമമാക്കണമെന്ന് സാമി സങ്കൽപിക്കുകയും പ്രതിഷ്ഠാക്കമ്മറ്റിക്കാർ അതിനെ അപ്രകാരം തന്നെ മനസ്സിലാക്കുകയും ചെയ്തു. പ്രതിഷ്ഠാക്കമ്മറ്റിയിലെ പ്രസിഡന്റ് ഡോക്ടർ പൽപ്പു മുതലായ യോഗ്യന്മാർ തന്നെ സ്വാമിയെ പ്രതിഷ്ഠ സംബന്ധിച്ച കാര്യങ്ങൾക്കായി മംഗലപുരംവരെ പോയി കൂട്ടിക്കൊണ്ടുവരികയും പ്രധാനമായ സകല ഏർപ്പാടുകളും ഒരുക്കങ്ങളും സ്വാമിയുടെ ഹിതവും കൽപനയും അനുസരിച്ച് ചെയ്കയും ചെയ്തു. പ്രതിഷ്ഠാ അവസരത്തിൽ ശിവഗിരിയിൽ സംഭാവനകളോടുകൂടി ഹാജരാകുന്നതിനും സ്വാമി പ്രതിഷ്ഠ നടത്തിയിട്ടുള്ള ചെറുതും വലുതും പൊതുവകയും പ്രത്യേകമാളുകളുടെ വകയുമായ 50-ഓളം ക്ഷേത്രങ്ങളിലേക്കു വിജ്ഞാപനങ്ങൾ അയക്കുകയും ചെയ്തപ്രകാരം പല ദൂരദേശത്തുള്ള ക്ഷേത്രങ്ങളിൽ നിന്നു വിഗ്രഹങ്ങളെ ആനപ്പുറത്ത് ഏഴുന്നള്ളിച്ച് ശിവഗിരിയിൽ എത്തുകയും ചെയ്തിരുന്നു. സ്വാമിതന്നെ നിശ്ചയിച്ച പ്രകാരം 1088 മേടം 18-ാം തിയതി രാത്രി 3 മണിക്ക് ശിവഗിരിയിൽ മഹാദേവ പ്രതിഷ്ഠയും 5 മണിക്ക് മദ്ധ്യേ ശാരദാ പ്രതിഷ്ഠയും സ്വാമി തന്റെ പാവനമായ കയ്യാൽതന്നെ നിർവഹിച്ചുകഴിഞ്ഞു. സംഗതികൾ എല്ലാം വിചാരിച്ചതിലും തുലോം അധികം ഭംഗിയായും കേമമായും മംഗളമായും കഴിഞ്ഞുകൂടി എന്ന് എങ്ങും ഒരുപോലെ ഉണ്ടായ ശ്രുതി സ്വാമിക്കും കമ്മറ്റിക്കും സമുദായത്തിനു മുഴുവനും ചാരിതാർത്ഥ്യജനകമായിരുന്നു. പ്രാധാന്യംകൊണ്ടും പ്രൗഡികൊണ്ടും ഇത്ര കേമമായ ഒരു മഹോത്സവം കേരളത്തിൽ തീയർ ഇതിനുമുൻപ് ഒരിക്കലും കൊണ്ടാടിയിട്ടില്ലെന്നുള്ളത് നിശ്ചയം തന്നെ.
	
	\newpage
	\chapter*{അദ്ധ്യായം പതിനഞ്ച്}
	
	ശിവഗിരിയിലെ പ്രതിഷ്ഠ കഴിഞ്ഞ് സ്വല്പ ദിവസം സ്വാമി ശിവഗിരിയിൽ താമസിച്ചു. അതിനു ശേഷം സന്യാസി ശിഷ്യന്മാരിൽ ചിലരെ മഠത്തിൽ താമസിപ്പാനും ചിലരെ പുറത്തുപോയി മതസംബന്ധമായ പ്രസംഗങ്ങൾ കൊണ്ടും മറ്റു പ്രകാരത്തിലും ജനങ്ങളുടെ ഗുണത്തിനായി യത്നിപ്പാനും ആജ്ഞാപിച്ച് അയച്ചുകൊണ്ട് സ്വാമി വീണ്ടും വടക്കോട്ടേക്ക് യാത്രചെയ്തു. വഴിക്ക് ചേർത്തല ഇറങ്ങുകയും സ്ഥലത്തെ സ്വജനങ്ങൾ സ്വാമിയെ ഭക്തി ബഹുമാനപൂർവ്വം സ്വീകരിച്ച് സൽക്കരിക്കയും, ആലുവായിൽ ഒരു മഠം പണിയുന്ന വകയ്ക്കായി കുറെ പണം ജനങ്ങൾ കാണിക്കവയ്ക്കയും ചെയ്തു. അവിടെനിന്നും സ്വാമി ആലുവായിലെത്തി താമസിച്ചു. ആ അവസരത്തിൽ തലശ്ശേരി ``ജഗന്നാഥ" ക്ഷേത്രത്തിൽ ഇടവമാസത്തെ ഇളംനീരഭിഷേകത്തിന് സ്വാമിയുടെ സാന്നിദ്ധ്യം ആവിശ്യപ്പെട്ട് ക്ഷേത്രഭാരവാഹികൾ വന്നു ക്ഷണിക്കയും വീണ്ടും സ്വാമി തലശ്ശേരിക്ക് പോകയും ചെയ്തു. ഈ യാത്രയിൽ സ്വാമി ആലുവായിൽ ഒരു സ്ഥാപനം ഉറപ്പിക്കുന്നതിനെപ്പറ്റി പ്രത്യക്ഷമായി ശ്രമിക്കുകയായിരുന്നു. ജഗന്നാഥക്ഷേത്രത്തിൽ അന്നത്തെ ഇളംനീരഭിക്ഷേകം സംബന്ധിച്ചുള്ള നടവരവിൽ ഒരു ഭാഗം ആലുവായിലെ സ്ഥാപനത്തിന്റെ ചിലവിലേക്കായി സമർപ്പിക്കാൻ ക്ഷേത്രഭാരവാഹികൾ തീർച്ചയാക്കുകയും അപ്രകാരം ചെയ്കയും ചെയ്തു. തലശ്ശേരിയിൽ നിന്നും സ്വാമി ഉടനെ ആലുവായ്ക്കു മടങ്ങുകയും അവിടെ മുൻപു സൂചിപ്പിച്ചിട്ടുള്ളതും, ഇപ്പോൾ സ്വാമി ആശ്രമം സ്ഥാപിച്ചിരിക്കുന്നതുമായ പുഴവക്കത്തുള്ള പറമ്പ് തീറെഴുതിവാങ്ങുകയും ചെയ്തു.
	
	കർക്കിടത്തിൽ വീണ്ടും സ്വാമി ശിവഗിരിയിലും അവിടെനിന്നും അരുവിപ്പുറത്തും എത്തി വിശ്രമിക്കയും സമീപപ്രദേശങ്ങളിൽ സഞ്ചരിക്കുകയും 1088 കന്നിമാസത്തോടുകൂടി ആലുവായ്ക്കു മടങ്ങുകയും ചെയ്തു. കുംഭമാസം ശിവരാത്രി സംബന്ധിച്ച് ആലുവായിൽ സ്വാമി ഒരു വലിയ സഭ വിളിച്ചുകൂട്ടുകയും ഒരു സംസ്കൃത വിദ്യാമന്ദിരം സ്ഥാപിക്ക മുതലായ സംഗതികളെപ്പറ്റി ആലോചിക്കയും ചെയ്തു. ഈ വിദ്യാമന്ദിരം സംബന്ധിച്ചു ജനങ്ങളെ പ്രേരിപ്പിക്കുന്നതിനും മറ്റുമായി വീണ്ടും സ്വാമി മലബാറിൽ സഞ്ചരിക്കുകയും അനേകം ധനവാന്മാർ അതിലേക്കായി വലിയ സംഖ്യകൾ കൊടുപ്പാൻ വാഗ്ദാനം ചെയ്കയും ചെയ്തു. മീനത്തിൽ സ്വാമി ശിവഗിരി, അരുവിപ്പുറം ഈ സ്ഥലങ്ങളിലും കൊല്ലത്തും ഏതാനും ദിവസം വിശ്രമിക്കയും മേടത്തിൽ ആലുവായ്ക്കു മടങ്ങി അവിടെനിന്നും സ്വല്പദിവസം നീലഗിരിയിൽ പോയി വിശ്രമിക്കയും മടക്കത്തിൽ പാലക്കാട്ട് ഇറങ്ങുകയും അവിടെ സ്വല്പം താമസിച്ച് ആലുവായിൽ തിരിയെ എത്തി വിശ്രമിക്കുകയും ചെയ്തു.
	
	1089-ാമാണ്ട് ചിങ്ങമാസം മുതൽ വൃശ്ചികംവരെ അധികദിവസവും സ്വാമി ആലുവായിൽ വിശ്രമിക്കയും അവിടെ നിന്നു വടക്കോട്ടു സഞ്ചരിക്കയും ചെയ്തുകൊണ്ടിരുന്നു. ഒടുവിൽ കൊല്ലം, കാർത്തികപ്പള്ളി ഈ സ്ഥലങ്ങളിൽ ക്ഷണിക്കപ്പെടുകയും കാർത്തികപ്പള്ളി ആലുംമൂട്ടിൽ തറവാട്ടിൽ നിന്നും മദ്രാസിൽ അവരുടെ വകയായുള്ള 13000ക. വിലപിടിക്കുന്ന ഒരു വീടും പറമ്പും ദാനമായി സ്വാമിക്ക് എഴുതിക്കൊടുക്കയും ചെയ്തു. ഇതിനിടയിൽ സ്വാമി ആലുവാപുഴവക്കിൽ തീറുവാങ്ങിയ പറമ്പിൽ ഒരു ചെറിയ ആശ്രമം പണികഴിപ്പിച്ചിട്ടുണ്ട്. ആശ്രമത്തിനടുത്തു പറവൂർ വടക്കേക്കര മൂത്തകുന്നത്തു ശ്രീനാരായണമംഗല ക്ഷേത്രക്കാർ വാങ്ങി സ്വാമിക്കു സമർപ്പിച്ചിരിക്കുന്ന, തീവണ്ടിസ്റ്റേഷനു സമീപമുള്ളതും സ്വാമി മുൻപ് വിഷൂചിക പിടിപെട്ടപ്പോൾ താമസിച്ചിരുന്നതുമായ കെട്ടിടത്തെ പുതുക്കി അതിൽ വിദ്യാർത്ഥികളും സന്യാസികളും മറ്റും താമസിക്കുന്നതിനുള്ള ഒരു മഠം ആക്കുകയും ഒരു അദ്ധ്യാപകനെ വച്ചു സംസ്കൃതം പഠിപ്പിക്കുന്നതിന് ഏർപ്പാടു ചെയ്യുകയും ചെയ്തിട്ടുണ്ട്. മകരത്തിൽ സ്വാമി ശിവഗിരിയിൽ ആദ്യമായി പൂയം മഹോത്സവത്തിന് ഏർപ്പാടു ചെയ്യുകയും ഉത്സവം മംഗളകരമായും കേമമായും നടക്കുകയും ചെയ്തു. ഈകൊല്ലം മേടമാസത്തോടടുത്ത്, അന്നും തിരുവിതാംകൂർ ചീഫ് ജസ്റ്റീസായിരുന്ന ഇപ്പോഴത്തെ ദിവാൻ മ.രാ.രാ. മന്ദത്തു കൃഷ്ണൻ നായർ അവർകൾ സ്വാമിയെ ആലുവായിൽവച്ചു സന്ദർശിച്ചിരുന്നു. മേടത്തിൽ ആലുവാവച്ച് അതികേമമായി നടത്തപ്പെട്ട എസ്. എൻ. ഡി. പി. യോഗത്തിന്റെ 11-ാം വാർഷിക പൊതുയോഗത്തിൽവച്ചു മാറിപ്പോയ ദിവാൻ രാജഗോപാലാചാരി അവർകൾ സ്വാമിയെ കാണാൻ മുൻകൂട്ടി പ്രതക്ഷിച്ചിരുന്നു എങ്കിലും ആ സന്ദർഭത്തിൽ സ്വാമി കുറ്റാലത്തു വിശ്രമിക്കയായിരുന്നാൽ സാധിച്ചില്ല. മേടം അവസാനത്തിൽ സ്വാമി മടങ്ങി ആലുവായ്ക്കു വരുന്ന മദ്ധ്യത്തിൽ കേരളിയ നായർസമാജം പ്രവർത്തകന്മാർ സ്വാമിയെ കോട്ടയത്തുവച്ചു നടത്തിയിരുന്ന ടി സമാജത്തിന്റെ വാർഷികയോഗത്തിൽ സംബന്ധിപ്പാനായി സൽക്കാരപൂർവ്വം കൂട്ടിക്കൊണ്ടുപോകയും സ്വാമി ടി സമാജത്തിൽ ഹാജരാകയും സഭയിൽ സന്നിഹിതരായിരുന്ന സർവ്വജനങ്ങളുടെയും അസാമാന്യമായ സ്നേഹബഹുമാനങ്ങൾക്ക് ഏക ലക്ഷ്യമായി തീരുകയും ചെയ്തു. കോട്ടയത്തുനിന്നു സ്വാമി ആലുവാ അദ്വൈതാശ്രമത്തിൽ എത്തി സ്വല്പദിവസം താമസിച്ചു. അവിടെനിന്നും, ഇടവത്തിൽ മലബാറിൽപോയി മടങ്ങിവന്നു. മിഥുനത്തിൽ കൊല്ലത്തുവന്നു. അവിടെ നിന്ന് ഏതാനും മാന്യഗൃഹസ്ഥന്മാരും ബ്രഹ്മചാരികളും ഒന്നിച്ച് കുറ്റാലം മുതലായ പല പുണ്യക്ഷേത്രങ്ങളും സന്ദർശിച്ചുകൊണ്ട് യാത്ര ചെയ്തു. മദ്രാസിലെത്തി സ്വല്പദിവസം താമസിക്കുകയും അവിടെനിന്ന് ബാംഗ്ലൂർവരെപോയി ആലുവായ്ക്കു മടങ്ങുകയും ചെയ്തു. 1090-ാമാണ്ടു ചിങ്ങമാസത്തിൽ സ്വാമി ചെങ്ങന്നൂർ, തിരുവല്ല ഈ താലൂക്കുകളിൽ സഞ്ചരിച്ചു. സ്വജനങ്ങൾ സ്വാമിയെ ഏറ്റവും ഭക്തിപുരസരം അതാതുസ്ഥലങ്ങളിൽ എതിരേൽക്കുകയും സൽക്കരിക്കുകയും ആലുവ സ്ഥാപിപ്പാൻ വിചാരിക്കുന്ന വിദ്യാലയത്തിന് ധനസഹായം ചെയ്യുകയും ചെയ്തു. ഈ യാത്രയിൽ അന്യവർഗങ്ങളിലെ പല മാന്യന്മാരും സ്വാമിയെ ആദരിക്കുകയും പൊതുവകയായുള്ള ചില സഭകളിൽ സ്വാമി സന്നിഹിതനാകുകയും ചെയ്തു. പുലയർ മുതലായ എളിയ വർഗ്ഗക്കാരുടെ മേൽ ജനങ്ങൾക്ക് അനുകമ്പ തോന്നേണ്ട ആവശ്യകതയെപ്പറ്റി ഈ സന്ദർഭത്തിൽ സ്വാമി പലരോടും സ്വകാര്യമായി ഉപദേശിക്കയും സ്വാമി പ്രത്യക്ഷമായി ആ വർഗ്ഗക്കാരോട് ഭേദമില്ലാതെയും പ്രത്യേകം സ്നേഹപൂർവ്വമായും പെരുമാറുകയും ചെയ്തു. സ്വാമിയുടെ ഈ സ്നേഹ പ്രകടനം ആ സ്ഥലത്തെ സ്ഥിതിക്ക് പ്രത്യേകം ആവശ്യം തന്നെ ആയിരുന്നു.
	
	ധനുമാസത്തിൽ സ്വാമി ശിവഗിരിയിലെത്തി സ്വല്പം വിശ്രമിക്കയും അവിടെനിന്നും അരുവിപ്പുറത്തു പോയി താമസിക്കയും സമീപസ്ഥലങ്ങളിൽ സഞ്ചരിച്ചുകൊണ്ടിരിക്കയും ചെയ്തു. നെയ്യാറ്റിൻകര പുലയർക്ക് ഈ അവസരത്തിൽ അന്യജാതിക്കാരിൽ നിന്നു നേരിട്ട ഉപദ്രവങ്ങളിൽ സ്വാമി അത്യന്തം സഹതപിക്കുകയും സ്വവർഗ്ഗക്കാരുടെ അനുകമ്പ അവരിൽ സ്ഥിരമായി പ്രവർത്തിക്കുന്നതിന് സ്വകാര്യമായി വേണ്ട ഉപദേശങ്ങൾ ചെയ്യുകയും ചെയ്തു. അനവധി പുലയരും അവരുടെ സമുദായ പ്രധാനികളും സ്വാമിയെവന്ന് സന്ദർശിച്ച് അനുഗ്രഹവും സദുപദേശങ്ങളും സ്വീകരിച്ചു. ഈ സന്ദർഭത്തിൽ സ്വാമിയുടെ ശ്രദ്ധയെ ആകർഷിച്ചിരുന്ന മറ്റൊരു വിഷയം എസ്. എൻ. ഡി. പി. യോഗത്തിന്റെ ആചാരസംബന്ധമായും മതസംബന്ധമായും മറ്റുമുള്ള ഉദ്ദേശങ്ങൾ നടപ്പിൽ വരുത്തുന്നതിനും അതിന്റെ പ്രവർത്തികളെ പൂർവ്വാധികം പ്രചാരപ്പെടുത്തുന്നതിനുമായി സ്വജനങ്ങൾ അധിവസിക്കുന്ന ദേശങ്ങൾ, അല്ലെങ്കിൽ കരകൾതോറും ദേശസഭകൾ ഏർപ്പെടുത്തുക എന്നുള്ളതായിരുന്നു. എസ്. എൻ. ഡി. പി. യോഗം പോലെതന്നെ ദേശസഭകളും കേരളത്തിന്റെ തെക്കേ അറ്റമായ നെയ്യാറ്റിൻകര താലൂക്കിൽ ആദ്യമായി സ്ഥാപിച്ചുതുടങ്ങുകയും, ക്രമേണ മറ്റു താലൂക്കുകളിലും രാജ്യങ്ങളിലും വ്യാപിപ്പിക്കയും ചെയ്യണമെന്നുള്ള വിചാരത്താൽ സ്വാമി അരുവിപ്പുറത്ത് അതിനായി കുറേദിവസം വിശ്രമിക്കയും സ്ഥലത്തെ ജനങ്ങളെ ആ വിഷയത്തിൽ പ്രേരിപ്പിക്കയും ചെയ്തുകൊണ്ടിരുന്നു. ഇതിനിടയിൽ തിരുവിതാംകൂറിന്റെ തെക്കുകിഴക്കേ അതിർത്തിയും തമിഴ് പ്രദേശവുമായ തോവാള, അഗസ്തീശ്വരം ഈ താലൂക്കുകളിലെ സ്വജനങ്ങൾ സ്വാമിയെ ക്ഷണിച്ച് ആ സ്ഥലങ്ങളിലേക്ക് കൊണ്ടുപോകയും തോവാള കടുക്കറ എന്ന സ്ഥലത്തും അഗസ്തീശ്വരത്ത് കോട്ടാർ നഗരത്തിലും സ്വാമി ഒന്നുരണ്ടു ദിവസം വിശ്രമിക്കയും ചെയ്തു. ഈ അവസരത്തിൽ സ്വാമി അവിടെ ആട്, കോഴി, മുതലായ ജന്തുക്കളെ ബലികഴിച്ചുവന്ന അനേകം ദുർദേവതകളുടെ പീഠങ്ങൾ ജനങ്ങളുടെ സമ്മതം വാങ്ങി ഇടിച്ചുകളയുകയും ചില ദേവിക്ഷേത്രങ്ങളിൽ നടന്നുവന്ന പ്രാണിഹിംസയെ നിർത്തൽചെയ്തു സാത്വികരീതിയിലുള്ള ആരാധനാക്രമം നടപ്പാക്കുകയും ചെയ്തു. ഈ നഗരത്തിൽ വാകയടിത്തെരുവു ആറുമുഖപ്പെരുമാൾ പിള്ളയാർ ദേവസ്വവും സമുദായവും വക കാര്യങ്ങൾക്കു സമുദായങ്ങളുടെ ഇടയിലുള്ള കക്ഷി മത്സരം നിമിത്തം ഉണ്ടായിരുന്ന കുഴപ്പങ്ങൾ എല്ലാം ഈ അവസരത്തിൽ സ്വാമി പറഞ്ഞു തീർത്തു രാജിപ്പെടുത്തുകയും മേലാൽ ഈ കുഴപ്പങ്ങൾ ഉണ്ടാകാതിരിപ്പാൻ ചില വ്യവസ്ഥകൾ ഏർപ്പെടുത്തുകയും ചെയ്തു. അവിടെനിന്നു സ്വാമി ശിവഗിരിയിൽ പൂയമഹോത്സവം സംബന്ധിച്ചു സന്നിഹിതനായിരിക്കേണമെന്നുള്ള ഉത്സവ ഭാരവാഹികളുടെ ഭക്തിപൂർവ്വകമായ അപേക്ഷ അനുസരിച്ച് ഉടനെ പുറപ്പെട്ടു ശിവഗിരിയിൽ എത്തി. ശിവഗിരി മഹാദേവ പ്രതിഷ്ഠ ശാരദാപ്രതിഷ്ഠയോടുകൂടി പെട്ടെന്നു ആലോചിച്ചു ചെയ്ത ഒരു സ്ഥാപനമാണെന്നും അതു വേണ്ടത്ര പൂർവ്വാലോചനയോടും ക്ഷേത്രം നിർമിക്ക മുതലായ ആവിശ്യങ്ങൾ പൂർത്തിയാക്കിയശേഷവും ചെയ്തതല്ലെന്നുമുള്ള വസ്തുത പരസ്യമാണ്. പ്രതിഷ്ഠാദിവസം മുതൽ ക്ഷേത്രത്തിന്റെ സ്ഥാനത്തിൽ ഒരു ചെറിയ ഒരു ഓലപ്പുര മാത്രമാണുള്ളത്. ക്ഷേത്രപ്പണി ജനങ്ങൾ വേഗം ആരംഭിച്ചു പൂർത്തിയാക്കുമെന്നായിരുന്നു സ്വാമി പ്രതീക്ഷിച്ചിരുന്നത്. ശിവഗിരി വിട്ടു പ്രതിഷ്ഠാനന്തരം സ്വാമി ആലുവായിൽ താമസം ആരംഭിച്ചതിനാൽ ജനങ്ങളുടെ ദൃഷ്ടി ക്രമേണ അങ്ങോട്ടേക്കു ആകർഷിക്കപ്പെടുകയും ശിവഗിരിയിലെ കാര്യങ്ങൾക്ക് ഉത്സാഹം ഏതാണ്ട് കുറഞ്ഞു തുടങ്ങുകയും ചെയ്തു. അതുകൊണ്ടുതന്നെയാണു ക്ഷേത്രപ്പണി ആരംഭിപ്പാൻ താമസം നേരിട്ടത്. എന്നാൽ സമുദായത്തിൽ പല യോഗ്യന്മാരും യോഗം ഭാരവാഹികളും സദാ അതിനെപ്പറ്റി ചിന്തിച്ചുകൊണ്ടും പണി ആരംഭിപ്പാൻ ഒരു നല്ല അവസരത്തെ പ്രതീക്ഷിച്ചുകൊണ്ടുമാണിരുന്നത്.
	
	ഉത്സവം സംബന്ധിച്ചു ശിവഗിരിയിൽ വിശ്രമിക്കുമ്പോൾ അവിടെ നിർമ്മിക്കേണ്ട ക്ഷേത്രത്തിന്റെ മാതൃകയേയും വലിപ്പത്തേയും മറ്റും പറ്റി സ്വാമി ഗാഢമായി ചിന്തിച്ചുകൊണ്ടിരിക്കയും ചില അഭിപ്രായങ്ങൾ എല്ലാം പ്രസ്താവിക്കയും ചെയ്തിട്ടുണ്ട്. വാക്കുകൊണ്ടും പ്രവർത്തികൊണ്ടും അനുമിക്കാവുന്നതായി സ്വാമിയുടെ അഗാധമായ ഹൃദയത്തിൽ ഇപ്പോൾ കിടക്കുന്ന മറ്റൊരു പാവനമായ അഭിപ്രായം തന്റെ ശിഷ്യന്മാരായ സന്യാസികളും ബ്രഹ്മചാരികളും അടങ്ങിയ ഒരു പ്രത്യേകസംഘം സ്ഥാപിച്ചും അതുമൂലം ജാതി മത ഭേതംകൂടാതെ പൊതുവിൽ നാട്ടിനും ജനങ്ങൾക്കും ഒന്നുപോലെ ആദ്ധ്യാത്മികമായ ശ്രേയസ്സും സദാചാരസംബന്ധമായും വിദ്യാഭ്യാസസംബന്ധമായും ഉള്ള അഭിവൃദ്ധിയും ഉണ്ടാകുന്നതിനു യത്നിപ്പാൻ വേണ്ട ഏർപ്പാടുകൾ ചെയ്യണമെന്നുള്ളതാകുന്നു.
	
	സ്വാമിയുടെ സംക്ഷിപ്തമായ ഈ ജീവചരിത്രം തൽകാലം ഇവിടെ അവസാനിപ്പിക്കേണ്ടിയിരിക്കുന്നു. ഇതിനെ വായനക്കാർ അതിമനോഹരമായ അവിടത്തെ ജീവചരിത്രഗാത്രത്തിന്റെ അസ്ഥികൂടം എന്നുമാത്രം കരുതിയാൽ മതി.
	
	\newpage
	\chapter*{{\large അനുബന്ധം}\\ 1090 മുതൽ മഹാസമാധി വരെ}
	
	ഒന്നിനോടും മമത പുലർത്താതെ അന്യരുടെ നന്മയ്ക്കു വേണ്ടി കർമ്മം ചെയ്യുക എന്ന സന്ദേശം ഗുരു സ്വന്തം ജീവിതം കൊണ്ടു ജനതയ്ക്കു കാണിച്ചു കൊടുത്തു.
	
	\begin{myverse}
		`അന്ത്യം വരെയും\\
		കർമ്മം ചെയ്തിങ്ങസംഗനായ്\\
		ചരിക്കുകയതല്ലാതി-\\
		ല്ലൊന്നൂം നരനു ചെയ്തിടാൻ'
	\end{myverse}
	
	എന്ന് ഈശാവാസ്യോപനിഷത്തിന്റെ വിവർത്തനത്തിലൂടെ അനുശാസിക്കുകയും ചെയ്തു. ഈ അനുശാസനം ഗുരു സ്വയം പാലിച്ചു എന്ന് ആശാൻ എഴുതിയ ജീവചരിത്രഭാഗത്തുനിന്നു തെളിയുന്നു. ആശാൻ പറഞ്ഞു നിറുത്തിയ വാക്കുകൾ ഉദ്ധരിക്കട്ടെ. ``വാക്കുകൊണ്ടും പ്രവൃത്തികൊണ്ടും അനുമിക്കാവുന്നതായി സ്വാമിയുടെ അഗാധമായ ഹൃദയത്തിൽ ഇപ്പോൾ കിടക്കുന്ന മറ്റൊരു പാവനമായ അഭിപ്രായം തന്റെ ശിഷ്യന്മാരായ സന്യാസികളും ബ്രഹ്മചാരികളും അടങ്ങിയ ഒരു പ്രത്യേക സംഘം സ്ഥാപിച്ചും അതുമൂലം ജാതിമതഭേതം കൂടാതെ പൊതുവിൽ നാട്ടിനും ജനങ്ങൾക്കും ഒന്നുപോലെ അദ്ധ്യാത്മികമായ ശ്രേയസ്സും സദാചാര സംബന്ധമായും വിദ്യാഭ്യാസ സംബന്ധമായും ഉള്ള അഭിവൃദ്ധിയും ഉണ്ടാകുന്നതിനു യത്നിപ്പാൻവേണ്ട ഏർപ്പാടുകൾ ചെയ്യണമെന്നുള്ളതാകുന്നു." അത്തരം കാര്യങ്ങൾ എത്രയും വേഗം ചെയ്തുതീർക്കണമെന്നു സ്വാമികൾ വിചാരിച്ചതായി അദ്ദേഹത്തിന്റെ പിൽക്കാലചര്യകളിൽ നിന്നൂഹിക്കാം. ആര്യസമാജ പ്രവർത്തകനായ സ്വാമി ശ്രദ്ധാനന്ദജിയുമായി നടന്ന സംഭാഷണത്തിനിടയിൽ ഗുരു പറഞ്ഞ നർമ്മങ്ങളിൽ ചെയ്തതൊന്നും പോരാ എന്നു ധ്വനിക്കുന്നു. പ്രസക്ത ഭാഗം താഴെ എഴുതുന്നു.
	
	ശ്രദ്ധാനന്ദ: അധഃകൃതരുടെ ഉദ്ധാരണത്തിനു സ്വാമികൾ പലതും ചെയ്തിട്ടുണ്ടെന്നറിയാം.
	
	ഗുരു: നാം ഒന്നും പ്രവർത്തിക്കുന്നില്ലല്ലൊ.
	
	ശ്രദ്ധാനന്ദ: പ്രവർത്തിയുണ്ടാകാതിരിക്കുകയില്ല. പ്രവർത്തിയുണ്ടെങ്കിലെ നിവൃത്തിയുള്ളുവല്ലോ.
	
	ഗുരു: ഇവിടെ നമുക്ക് ഒരു നിവൃത്തിയും ഇല്ല.
	
	ഇങ്ങനെ ചെയ്തില്ലെന്നും ചെയ്യുന്നില്ലെന്നും തോന്നിയതിനാൽ സ്വാമികൾ അവിരതം പ്രവർത്തനനിരതനായി. നിറുത്താത്ത സഞ്ചാരമായിരുന്നു പ്രവർത്തനങ്ങളിൽ മുഖ്യ ഇനം. സഞ്ചാരത്തിനിടയിൽ സന്മനസ്സുള്ളവരിൽ നിന്നു പൊതുസ്ഥാപനങ്ങൾ നിർമ്മിക്കാൻ സംഭാവനകൾ സ്വീകരിക്കുക പതിവായിരുന്നു. ആരാധനാലയം, ആശ്രമം, പള്ളിക്കൂടം, തൊഴിൽശാല എന്നിവ സ്ഥാപിക്കാൻ ഈ സംഭാവനകൾ പ്രയോജനപ്പെടുത്തി. സാധാരണക്കാരുടെ അജ്ഞതയും ദാരിദ്ര്യവും പരമാവധി ദൂരികരിക്കുക എന്നതായിരുന്നു ആത്യന്തിക ലക്ഷ്യം.
	
	പ്രകൃതിരമണീയകത്വം വഴിഞ്ഞൊഴുകുന്ന പ്രദേശങ്ങൾ മനസ്സിന്റെ എകാഗ്രത മൂലം നേടാവുന്ന ശാന്തിക്കുതകുമെന്നു സ്വാമികൾ അനുഭവിച്ചറിഞ്ഞിരുന്നു. ആ നിലക്കു പെരിയാറിന്റെ തീരം സ്വാമികളെ നേരത്തേ തന്നെ ആകർഷിച്ചു കാണണം. ശങ്കരാചാര്യന്റെ ജനന സ്ഥലമായ കാലടിക്കടുത്താണെന്ന കാര്യവും സ്വാമികൾ ഓർത്തിരിക്കണം. മാത്രമല്ല കൊല്ലം തോറും ശിവരാത്രിക്ക് ആയിരക്കണക്കിനാളുകൾ ബലിയിടാൻ ആലുവാമണപ്പുറത്തു വന്നുചേർന്നു പിരിഞ്ഞു പോകാറുണ്ട്. മലയിടുക്കിൽക്കൂടി ഒരു അരുവി കുത്തിച്ചാടുന്നതു കാണുന്ന എൻജിനിയറുടെ മനസ്സിൽ അവിടെ ഒരണ കെട്ടിയാലോ എന്ന ചിന്ത ഉയരും. ഈ ജനക്കൂട്ടത്തെ കണ്ട സ്വാമികളുടെ മനസ്സിലും സമാനമായ ഒരു വിചാരം ഉയർന്നിരിക്കണം കൊല്ലവർഷം 1090-ൽ ആലുവാ പുഴയുടെ തീരത്ത് അദ്വൈതാശ്രമം സ്ഥാപിക്കാനുള്ള ശ്രമം ആരംഭിച്ചു. അവിടെ സ്ഥലം വാങ്ങാൻ വേണ്ട ധനം സംഭരിക്കാൻ സ്വാമികൾ തന്നെ ഭക്തന്മാരെ സന്ദർശിച്ചു. പണം തികഞ്ഞപ്പോൾ സ്ഥലം വാങ്ങി ആദ്യം അവിടെ ഒരു പർണശാലയും മഠവും പണി കഴിപ്പിച്ചു. അതിനുശേഷം അതിനോടു ചേർന്നു ഒരു വിദ്യാലയം സ്ഥാപിക്കാനുള്ള പ്രയത്നം തുടങ്ങി.
	
	സ്വാമികൾ അവധൂതനായി നടന്നകാലത്തെ സഞ്ചാരശീലം ഒരിക്കലും ഉപേക്ഷിച്ചില്ല. ``ചരൈവേതി ചരൈവേതി" സ്വാമികളെക്കുറിച്ചു തികച്ചും ശരിയായിരുന്നു. അക്കൊല്ലത്തെ സഞ്ചാരത്തിനിടയിൽ പഴനി മല സന്ദർശിച്ചു.
	
	ഏതെങ്കിലും പ്രദേശത്തുള്ള ആളുകൾ ചേർന്ന് അവർക്കൊരു ആരാധനാലയം വേണമെന്ന് അപേക്ഷിക്കുകയും സ്വാമികൾക്കത് ബോദ്ധ്യപ്പെടുകയും ചെയ്താൽ സ്ഥാപിച്ചു കൊടുക്കുക പതിവുണ്ട്. അനാഡംബരമായ ആരാധനാലയം ശുചിത്വം, ഈശ്വരഭക്തി, സംഘടനാബോധം, മിതവ്യയശീലം, വിദ്യാഭ്യാസം, ജാതിമതഭേദനിരാസം എന്നിവയ്ക്ക് ഉതകുമെന്നു സ്വാമികൾ കരുതിയിരുന്നു. അക്കൊല്ലം കമ്മാര സമുദായക്കാർക്കുവേണ്ടി സിദ്ധേശ്വരം ക്ഷേത്രം സ്ഥാപിച്ചു. അഞ്ചുതെങ്ങിൽ ജ്ഞാനേശ്വരക്ഷേത്രത്തിലെ പ്രതിഷ്ഠ നടത്തിയതും അതേ ആണ്ടിൽ ആണ്.
	
	ജാതിയും അതിന്റെ പേരിലുള്ള അനാചാരങ്ങളും മനുഷ്യമനസ്സിനെ പ്രകടമായി മലിനമാക്കിയിരുന്ന കാലമാണല്ലൊ അത്. ഈ മാലിന്യം തുടച്ചു മാറ്റാൻ ഗുരു സ്വകീയമായ മാർഗ്ഗങ്ങളിലൂടെ പരിശ്രമിച്ചുകൊണ്ടിരുന്നു. അക്കൊല്ലം അതിന്റെ ബീഭത്സരൂപം തലനീട്ടിയ ഒരു സംഭവം ഉണ്ടായി. ഊരൂട്ടമ്പലം പെൺപള്ളീക്കൂടത്തിൽ പുലയപ്പെൺകുട്ടികളെ ചേർക്കാനുള്ള ശ്രമത്തെ ചില നായന്മാർ എതിർത്തു. കലഹവും അടിപിടിയും ചെറിയ തോതിൽ ലഹളയും നടന്നു. പുലയരെ ഉപദ്രവിക്കാൻ ചില ഈഴവ പ്രമാണികളും നായന്മാരുടെ കൂട്ടുചേർന്നു. സ്വാമികൾ സ്ഥലം സന്ദർശിച്ചു. മർദ്ദിതർക്ക് ആശ്വാസം അരുളി, മനുഷ്യത്വത്തിനേറ്റ ഈ മുറിവു ഗുരുവിനെ ഏറെ വേദനിപ്പിച്ചു.
	
	കേവലം മനുഷ്യകൃതമായ ജാതിഭേദത്തിന്റെ കെടുതികൾ ഇല്ലാതാക്കാൻ സ്വാമികൾ ഇടതടവില്ലാതെ പ്രവർത്തിച്ചുകൊണ്ടിരുന്നു. ഗുരു അക്കൊല്ലം ചെങ്ങന്നൂർ, തിരുവല്ല, ചങ്ങനാശ്ശേരി തുടങ്ങിയ സ്ഥലങ്ങൾ സന്ദർശിക്കുകയുണ്ടായി. അവിടെ പുലയരുടെ ഒരു യോഗം വിളിച്ചുകൂട്ടി. കൂട്ടത്തിൽ കഴിവുള്ള ഒരു കുട്ടിയെകൊണ്ട് പ്രസംഗം ചെയ്യിച്ചു. ആ കുട്ടിയെ സ്വാമികൾ കൂടെ കൊണ്ടുപോരികയും ചെയ്തു. ജാതിചിന്ത വേരറുത്തുകളയാൻ ഏറ്റവും പറ്റിയ മാർഗം എന്ന നിലയിൽ പുലയ വിദ്യാർത്ഥികളെ ചേർത്തു മിശ്രഭോജനം നടത്തി.
	
	ആലുവായിൽ സംസ്കൃത പാഠശാല സ്ഥാപിക്കാനുള്ള ശ്രമം 1091-ൽ ഫലിച്ചു. മനുഷ്യൻ നന്നാകാനുള്ള ഏറ്റവും പറ്റിയ ഉപാധി അവന് അറിവുണ്ടാക്കിക്കൊടുക്കുകയാണെന്നു ഗുരു കണ്ടിരുന്നു.
	
	ആ വർഷം തന്നെ രമണമഹർഷിയുടെ ക്ഷണം സ്വീകരിച്ചു ഗുരു തിരുവണ്ണാമലയിൽ പോയി വിശ്രമിച്ചു. ആ അവസരത്തിൽ എഴുതിയതാണു മുനിചര്യാപഞ്ചകം. രണ്ടു യോഗീശ്വരന്മാർ മൗനത്തിലൂടെ എങ്ങനെ പരസ്പരം അറിഞ്ഞു എന്നതിനു നിദർശനമായിരുന്നുവത്രേ അവരുടെ കൂടിക്കാഴ്ച. ഒരുദിവസം ഗുരുവിന്റെ മുമ്പിൽ കുറെ നേരം ഇരുന്നശേഷം മഹർഷി ആശ്രമത്തിന്റെ ഉള്ളിലേയ്ക്കുപോയി. സ്വാമികൾ രമണ മഹർഷിയുടെ അനുയായികളോടു ചോദിച്ചു: ``മനസ്സിലായോ?" മനസ്സിലായി എന്നവർ മറുപടി പറഞ്ഞപ്പോൾ ഗുരു സ്വന്തം അനുചരന്മാരോട് അതേ ചോദ്യം ആവർത്തിച്ചു. അവരും അതേ ഉത്തരം പറഞ്ഞപ്പോൾ ഗുരു മന്ദഹാസപൂർവ്വം മൊഴിഞ്ഞു: ``ഓഹോ, അപ്പോൾ നമുക്കു മാത്രമേ മനസ്സിലാവാതുള്ളു."
	
	മനുഷ്യർക്ക് മനുഷ്യത്വം എന്ന ജാതിയല്ലാതെ വേറെ ജാതിയില്ലെന്നു ഗുരു എപ്പോഴും ഓർമ്മിപ്പിക്കുമായിരുന്നു. 1091-ൽ ജാതി ഇല്ല എന്നൊരു വിളംബരം സ്വാമികൾ പുറപ്പെടുവിച്ചു. 1091 ഇടവത്തിലെ പ്രബുദ്ധ കേരളത്തിൽ വന്ന ആ വിളംബരം ഇങ്ങനെയായിരുന്നു. ``നാം ജാതിഭേദം വിട്ടിട്ട് ഇപ്പോൾ ഏതാനും സംവത്സരങ്ങൾ കഴിഞ്ഞിരുക്കുന്നു. എന്നിട്ടും ചില പ്രത്യേക വർഗ്ഗക്കാർ നമ്മെ അവരുടെ കൂട്ടത്തിൽ പെട്ടതായി വിചാരിച്ചും പ്രവർത്തിച്ചും വരുന്നതായും അതു ഹേതുവാൽ പലർക്കും നമ്മുടെ വാസ്തവത്തിനു വിരുദ്ധമായ ധാരണയ്ക്കിടവന്നിട്ടുണ്ടെന്നും അറിയുന്നു. നാം പ്രത്യേക ജാതിയിലോ മതത്തിലോ ഉൾപ്പെടുന്നില്ല. വിശേഷിച്ചും നമ്മുടെ ശിഷ്യവർഗത്തിൽ നിന്നും മേൽ പ്രകാരമുള്ളവരെ മാത്രമേ നമ്മുടെ പിൻഗാമിയായി വരത്തക്കവണ്ണം ആലുവാ അദ്വൈതാശ്രമത്തിൽ ശിഷ്യസംഘത്തിൽ ചേർത്തിട്ടുള്ളൂ എന്നും മേലും ചേർക്കയുള്ളൂ എന്നും വ്യവസ്ഥപ്പെടുത്തിയിരിക്കുന്നതുമാകുന്നു."
	
	1092-ൽ സ്വാമികളുടെ ഷഷ്ടിപൂർത്തി കേരളത്തിൽ അങ്ങോളമിങ്ങോളം ആഘോഷിച്ചു. കേരളത്തിനു പുറത്തു സിലോൺ, സിങ്കപ്പൂർ, മലയ എന്നിവിടങ്ങളിലും ആഘോഷം അല തല്ലി. സ്വാമി അന്നും നിർലിപ്തനായി നിലകൊണ്ടു. തിരുവനന്തപുരത്തു കൈതമുക്കിൽ ഷഷ്ടിപൂർത്തി സ്മാരകമായി ഒരു മന്ദിരം പണികഴിപ്പിച്ചു. പണിയുടെ ചുമതല കുമാരനാശാന് ആയിരുന്നു. കെട്ടിടത്തിന്റെ ഉദ്ഘാടന കർമ്മം അർഹിക്കുന്ന ഗൗരവത്തോടെ നടന്നു. അന്നു തിരുവിതാംകൂർ ദിവനായിരുന്ന ദിവാൻ ബഹദൂർ എം. കൃഷ്ണൻ നായരായിരുന്നു ഉദ്ഘാടകൻ.
	
	മിശ്രവിവാഹവും മിശ്രഭോജനവും ജാതി നശീകരണത്തിന് ഫലപ്രദമാണെന്നു സ്വാമികൾ കണ്ടിരുന്നു. ചെറായിയിൽ മിശ്രഭോജനം നടത്തി സഹോദരൻ അയ്യപ്പൻ പുലയനയ്യപ്പൻ എന്ന പേരു സമ്പാദിച്ചത് അക്കൊല്ലമാണ്. യാഥാസ്ഥിതികരുടെ അർഥശൂന്യമായ എതിർപ്പിന്റെ ശക്തി കുറയ്ക്കാൻ സ്വാമികളുടെ ഒരു സന്ദേശം കിട്ടിയാൽ കൊള്ളാമെന്നു സഹോദരൻ ആഗ്രഹിച്ചു. അങ്ങനെ എഴുതിക്കൊടുത്തതാണ് താഴെ വരുന്ന പ്രഖ്യാതമായ സന്ദേശം: ``മനുഷ്യരുടെ മതം, വേഷം, ഭാഷ മുതലായവ എങ്ങിനെയിരുന്നാലും അവരുടെ ജാതി ഒന്നായതുകൊണ്ടു അന്യോന്യം വിവാഹവും പന്തിഭോജനവും ചെയ്യുന്നതിനു യാതൊരു ദോഷവും ഇല്ല."
	
	അക്കൊല്ലം സ്വാമികൾ മദ്രാസും കാഞ്ചീപുരവും സന്ദർശിച്ചു. ചിന്താദ്രിപ്പേട്ടിലും കാഞ്ചീപുരത്തും അദ്വൈതാശ്രമങ്ങൾ സ്ഥാപിച്ചു. അവിടെ നിന്നു സ്വാമികൾ ബാഗ്ലൂർ പോയി ഡോ. പല്പുവിനെ കണ്ടു. കൂർക്കഞ്ചേരിയിൽ ക്ഷേത്ര പ്രതിഷ്ഠ നടത്തിയതും അതേ വർഷം തന്നെയാണ്.
	
	സ്വാമികളുടെ സഞ്ചാര പ്രിയത്വവും പതിതോദ്ധാരണ തത്പരതയും സിലോൺ (ശ്രീലങ്ക) നിവാസികൾക്കും അനുഗ്രഹമായിട്ടുണ്ട്. രണ്ടു പ്രാവിശ്യം സ്വാമികൾ സിലോൺ സന്ദർശിച്ചു. ആദ്യത്തെ യാത്ര 1094-ൽ ആയിരുന്നു. ആ യാത്രയിലാണു സ്വാമികൾ ഒന്നാമതായി കാവിവസ്ത്രം ധരിച്ചത്. അതുവരെ ശുഭ്രവസ്ത്രം ധരിക്കുകയായിരുന്നു പതിവ്. രാമേശ്വരത്തു വച്ചു ശിഷ്യന്മാരുടെ ഇംഗിതമനുസരിച്ചു സ്വാമികൾ കാവി കൈക്കൊണ്ടു. ആചാരങ്ങൾ ഒന്നും ഉണ്ടായില്ല. ആ തവണ ഗുരു പന്ത്രണ്ടു ദിവസം സിലോണിൽ താമസിച്ചു.
	
	ഈഴവരുടെ ഇടയിൽ പല അവാന്തര ജാതികളും ഉണ്ടായിരുന്നു. അവരിൽ ചിലരെ വേറേ ചിലർ താഴ്ന്നവരായി ഗണിച്ചുപോന്നു. അത്തരത്തിൽപ്പെട്ടവരായിരുന്നു പിച്ചിനാടു കുറുപ്പ് അല്ലെങ്കിൽ കണിക്കുറുപ്പ് എന്നറിയപ്പെട്ടിരുന്ന ഗണകവർഗം. അവരുടെ ഇടയിൽ കൃഷ്ണൻ വൈദ്യൻ എന്നൊരാൾക്കു വൈദ്യവൃത്തിമൂലം സാമാന്യം ഉയർന്ന നിലയുണ്ടായിരുന്നു. എന്നാൽ ജാതിയുടെ പേരിൽ മാന്യത കിട്ടിയില്ല. സ്വാമികളുടെ അദ്ധ്യക്ഷതയിൽ യോഗം കൂടി അദ്ദേഹത്തെ ഈഴവനായി പ്രഖ്യാപിച്ചു. യോഗത്തിൽ ഈഴവരും നായന്മാരും ക്രിസ്ത്യാനികളും സംബന്ധിച്ചിരുന്നു. യോഗാനന്തരം നടന്ന പൊതു സദ്യയിൽ എല്ലാവരും പങ്കുകൊണ്ടു. അങ്ങനെ ഇരിക്കെ ഒരിക്കൽ ഈഴവരുടെ പ്രതിനിധിയായി കൃഷണൻ വൈദ്യനു നാമനിർദ്ദേശം കിട്ടി. ഈഴവർ മുറുമുറുത്തു. അന്നു പേഷ്കാരായിരുന്ന ഉള്ളൂർ സ്വാമികളുടെ സാക്ഷ്യപത്രം ഉണ്ടായാൽ മതിയെന്നു പറഞ്ഞു. കൃഷ്ണൻ വൈദ്യൻ പരിശുദ്ധ ഈഴവനാണെന്നു ഗുരു സാക്ഷ്യപത്രം കൊടുത്തു ``പരിശുദ്ധ" എന്ന വിശേഷണത്തിനു ചെത്താത്ത എന്ന വ്യാഖ്യാനം മുറുമുറുപ്പുകാർക്കുള്ള ഒരടിയായിരുന്നു. 1094-ൽ ആണിതു നടന്നത്.
	
	സ്വാമികൾ നടത്തിയ പ്രതിഷ്ഠകളിൽ ചിന്തിക്കാൻ വക നൽകുന്ന തരം വൈവിദ്ധ്യം സംദൃശ്യമാണ്. നെയ്യാറിൽ നിന്നു മുങ്ങിയെടുത്ത അമ്മിക്കുഴവി പോലുള്ള കല്ലാണു ഗുരു ശിവലിംഗം എന്ന ഭാവനയിൽ അരുവിപ്പുറത്തു പ്രതിഷ്ഠിച്ചത്. അതിന്റെ ശിവകരത്വം-ശങ്കരത്വം-മംഗളകരത്വം-ഇന്നും കേരളീയർ അനുഭവിക്കുന്നുണ്ടല്ലൊ. പിന്നീടു പല ക്ഷേത്രങ്ങളിലും സുബ്രഹ്മണ്യൻ, ദേവി എന്നിവരുടെ വിഗ്രഹങ്ങൾ പ്രതിഷ്ഠിച്ചു. ഒരു നയിനാർ ക്ഷേത്രം പുതുക്കിപ്പണിതു. മാടൻ, മറുത മുതലായ ദേവതകളെ മാറ്റി മറ്റു ദേവന്മാരെ പ്രതിഷ്ഠിച്ചു. 1095-ൽ കാരമുക്ക് എന്ന സ്ഥലത്തു ഗുരു ക്ഷേത്രം സ്ഥാപിച്ചപ്പോൾ അവിടെ ഒരു ദീപം മാത്രമാണ് പ്രതിഷ്ഠിച്ചത്. വെളിച്ചം അറിവിന്റെ പ്രതീകമാണ്, അറിവു തന്നെ പരം പൊരുൾ.
	
	1096-ൽ സമസ്ത കേരള സഹോദര സമ്മേളനം ആലുവായിൽ നടന്നു. ഒരു ജാതി, ഒരു മതം, ഒരു ദൈവം, മനുഷ്യന് എന്ന അതിസരളവും അത്യന്തനൂതനവും ആയ മഹാസന്ദേശം പുറപ്പെടുവിച്ചത് ആ വർഷം ആണ്.
	
	മദ്യവർജനം സ്വാമികളുടെ പ്രവർത്തനങ്ങളുടെ ഒരു ഭാഗമാണ്. മദ്യം വിഷമാണ്, അതുണ്ടാക്കരുത്, കൊടുക്കരുത്, കുടിക്കരുത് എന്നും ചെത്തുകാരന്റെ ദേഹം നാറും, വസ്ത്രം നാറും, അവൻ തൊട്ടതൊക്കെ നാറും എന്നും സ്വാമികൾ 1096-ൽ പ്രഖ്യാപിച്ചു.
	
	അന്നു കേരളത്തിൽ ഈഴവരുടെ ഇടയിൽ ഒരു മത പരിവർത്തന കോലാഹലം നടന്നിരുന്നു. ഈഴവർ ഒന്നടങ്കം ക്രിസ്തുമതത്തിലേക്കു മാർഗം കൂടിയാൽ ജാതീയമായ അവശതകൾ മുഴുവൻ ഒറ്റയടിക്കു തീരുമെന്ന് ഒരു കൂട്ടർ വാദിച്ചു. മതം മാറുന്നെങ്കിൽ അതു ബുദ്ധമതത്തിലേക്കായിരിക്കണമെന്നു ചിലർ അഭിപ്രായപ്പെട്ടു. ഇസ്ലാം മതത്തിലെ സാഹോദര്യഭാവന വേറെ ചിലരെ ആകർഷിച്ചു. വാദപ്രതിവാദം കൊടുമ്പിരിക്കൊണ്ട അവസരത്തിൽ സ്വാമികൾ തമിഴുനാട്ടിൽ ആയിരുന്നു. തിരിച്ചു വന്നപ്പോൾ സഹോദരനയ്യപ്പനുമായി മതത്തെക്കുറിച്ചൊരു ചർച്ച നടത്തി. അതിൽ നിന്ന് ഉരുത്തിരിഞ്ഞു വന്നതാണ് ``മത മേതായാലും മനുഷ്യൻ നന്നായാൽ മതി" എന്ന മഹാസന്ദേശം.
	
	കാരമുക്കിലെ ദീപ മാത്ര പ്രതിഷ്ഠയിൽ പൊതുജനത്തിനു സംതൃപ്തി പോരെന്നു സ്വാമികൾക്കു തോന്നിയിരിക്കണം. 1097-ൽ മുരുക്കുംപുഴയിൽ ദീപത്തിനു പുറമെ ``ഓം സത്യം, ധർമ്മം, ദയ, ശാന്തി" എന്നെഴുതിയതും പ്രതിഷ്ഠിച്ചു. ഇന്ന് ആ ക്ഷേത്രത്തിൽ പല ദേവതകൾക്കും ഉപപ്രതിഷ്ഠകളായി സ്ഥാനം നൽകിയിരിക്കുന്നു. കൊല്ലം പെരിനാട് എസ്. എൻ. ഡി. പി.യുടെ അക്കൊല്ലത്തെ യോഗത്തിൽ സ്വാമികൾ സന്നിഹിതരായി.
	
	1098 ശിവഗിരിയുടെ ചരിത്രത്തിൽ വളരെ പ്രധാനപ്പെട്ട വർഷമാണ്. അക്കൊല്ലം രവീന്ദ്രനാഥ ടാഗോർ സി. എഫ്. ആൻഡ്രൂസിനോടൊപ്പം സ്വാമികളെ സന്ദർശിച്ചു. സ്വാമികളുടെ അഭൗമമായ പ്രഭാവം മഹാകവിയെ വളരെ ആകർഷിച്ചു.
	
	``കേരളത്തിലെ സ്വാമി നാരായണഗുരുവിനേക്കാൾ ആത്മീയമായി മാഹാത്മ്യമുള്ള ആരെയും ഞാൻ കണ്ടുമുട്ടിയിട്ടില്ല. എന്തിന് ആത്മീയസിദ്ധിയിൽ അദ്ദേഹത്തോളമെങ്കിലും മഹാനായ ആരെയും കണ്ടിട്ടില്ല. ദിവ്യപ്രഭയുടെ സ്വയം പ്രകാശമാനമായ വെളിച്ചത്തിൽ തിളങ്ങുന്ന ആ ജ്വലിക്കുന്ന മുഖവും വിദൂര ചക്രവാളത്തിൽ അങ്ങകലെയുള്ള ഒരു ബിന്ദുവിൽ നോട്ടം ഉറപ്പിച്ച ആ യോഗനയനങ്ങളും ഞാൻ ഒരിക്കലും മറക്കുകയില്ലെന്ന് എനിക്കുറപ്പുണ്ട്." എന്നു അന്ന് ടാഗോർ രേഖപ്പെടുത്തി.
	
	1098-ൽ പാണാവള്ളി ശ്രീകണ്ഠേശ്വര പ്രതിഷ്ഠ നടത്തി.
	
	പിറ്റേക്കൊല്ലം വ്യസനകരമായ ഒരു സംഭവം ഉണ്ടായി. പല്ലന റെഡീമർ ബോട്ടപകടത്തിൽ കുമാരനാശാൻ മരിച്ചു. അന്നു വെളുപ്പാൻ കാലത്തു സ്വാമികൾ ശിവഗിരിയിൽ വിശേഷാൽ പ്രാർത്ഥന നടത്തിച്ചു. ആശാന്റെ ചരമവാർത്ത പിന്നീടാണ് അറിഞ്ഞത്.
	
	വൈക്കത്തെ ക്ഷേത്രത്തിന്റെ പരിസരത്തുള്ള റോഡുകളിൽ കൂടി അവർണർക്ക് സഞ്ചരിക്കാനുള്ള സ്വാതന്ത്ര്യം നേടിയെടുക്കുന്നതിനുവേണ്ടി 1099-ൽ വൈക്കം സത്യാഗ്രഹം ആരംഭിച്ചു. ഗാന്ധിജിയുടെ അനുവാദത്തോടെ ടി. കെ. മാധവനാണ് ഇതിനു നേതൃത്വം നൽകിയത്. സത്യാഗ്രഹികൽക്കുവേണ്ടി സ്വാമികൾ വെല്ലൂർ മഠം വിട്ടുകൊടുത്തു. ഒരു സത്യാഗ്രഹഫണ്ട് ആരംഭിച്ചു. സ്വാമികൾ സത്യാഗ്രഹാശ്രമം സന്ദർശിച്ചു. അവിടുത്തെ എർപ്പാടുകളെപ്പറ്റി സംതൃപ്തി പ്രകടിപ്പിച്ചു.
	
	അക്കൊല്ലം സ്വാമികൾ ആലുവായിൽ സർവ്വമത സമ്മേളനം വിളിച്ചുകൂട്ടി. വാദിക്കാനും ജയിക്കാനുമല്ല, അറിയാനും അറിയിക്കാനുമാണ് എന്ന് സ്വാമികൾ സമ്മേളന കവാടത്തിൽ എഴുതിവപ്പിച്ചു. ആശയപരമായ ഏതു സംവാദവും ഈ മനോഭാവത്തോടെയാണു നടത്തേണ്ടതെന്ന് എടുത്തുപറയേണ്ടതില്ലല്ലൊ. രണ്ടു ദിവസമായി നടന്ന സമ്മേളനത്തിൽ അനേകം മതപണ്ഡിതന്മാർ പങ്കെടുത്തു. സ്വാമികൾ രണ്ടു ദിവസവും സമ്മേളനവേദിയിൽ ഉണ്ടായിരുന്നു. സമ്മേളനത്തിന്റെ ഉദ്ദേശലക്ഷ്യങ്ങൾ വിവരിച്ചുകൊണ്ട് സത്യവ്രത സ്വാമികൾ സ്വാഗത പ്രസംഗം നടത്തി. സി.വി. കുഞ്ഞുരാമൻ കൃതജ്ഞത പ്രകാശിപ്പിച്ചു. സിലോണിൽ നിന്നൊരു ബുദ്ധഭിക്ഷു സമ്മേളനത്തിൽ സംബന്ധിച്ചിരുന്നു. ബുദ്ധമതത്തെ കുറിച്ചു സംസാരിച്ചത് മഞ്ചേരി രാമയ്യരും രാമകൃഷ്ണയ്യരും ആണ്. ക്രിസ്തുമതത്തെക്കുറിച്ചു കെ. കെ. കുരുവിളയും ഇസ്ലാം മതത്തെക്കുറിച്ചു മഹമ്മദു മൗലവിയും ബ്രഹ്മസമാജത്തെക്കുറിച്ചു സ്വാമി ശിവപ്രസാദും ആര്യസമാജത്തെക്കുറിച്ചു പണ്ഡിറ്റ് ഋഷിറാമും പ്രസംഗിച്ചു. സമ്മേളനാനന്തരം ഒരു ഗ്രൂപ്‌ഫോട്ടോ എടുക്കുകയും സ്വാമികൾ ഒരു സന്ദേശം പുറപ്പെടുവിക്കുകയും ചെയ്തു: ``എല്ലാ മതങ്ങളുടേയും പരമോദ്ദേശ്യം ഒന്നാണെന്നും ഭിന്നമതാനുയായികൾ തമ്മിൽ കലഹിച്ചിട്ടാവിശ്യമില്ലെന്നും ഈ മഹാ സമ്മേളനത്തിൽ നടന്ന പ്രസംഗങ്ങൾ വെളിപ്പെടുത്തിയിരിക്കുന്നതിനാൽ നാം ശിവഗിരിയിൽ സ്ഥാപിക്കാൻ ഉദ്ദേശിക്കുന്ന പാഠശാലയിൽ എല്ലാ മതങ്ങളും പഠിക്കുന്നതിനുവേണ്ട എല്ലാ സൗകര്യങ്ങളും കൂടി ഉണ്ടാകണമെന്നു വിചാരിക്കുന്നു. ഈ സ്ഥാപനത്തിന്റെ നടത്തിപ്പിന് അഞ്ചു ലക്ഷം രൂപ പൊതുജനങ്ങളിൽ നിന്നും ലഭിക്കുന്നതിന് എല്ലാവരും സഹായിക്കുമെന്ന് വിചാരിക്കുന്നു." ഇതായിരുന്നു സ്വാമികളുടെ സന്ദേശം.
	
	1100 കുംഭത്തിൽ ഗാന്ധിജി ശിവഗിരി സന്ദർശിച്ചു. ജാതി, അയിത്തം മുതലായവയ്ക്കെതിരായി ഗുരു നടത്തുന്ന പരിശ്രമങ്ങൾ ഗാന്ധിജിയെ സന്തുഷ്ടനാക്കി. ഹരിജനക്കുട്ടികൾ പ്രാർത്ഥന ചൊല്ലിയതും അവരെ ശാന്തിക്കാരും പാചകക്കാരും മറ്റും ആക്കിയതും അത്യന്തം ആഹ്ലാദത്തോടെയാണ് ഗാന്ധിജി വീക്ഷിച്ചത്. ശിവഗിരിയിലും തിരുവനന്തപുരത്തും നടന്ന യോഗങ്ങളിൽ ഗാന്ധിജി ആ സന്ദർശനത്തിലുള്ള ആനന്ദം വെളിവാക്കി. സ്വാമികളെ കണ്ടത് മഹാഭാഗ്യമായി കരുതുന്നു എന്നു ഗാന്ധിജി പറഞ്ഞു. വൈക്കം സത്യാഗ്രഹത്തോടുള്ള സ്വാമിയുടെ സമീപനവും ഗാന്ധിജിക്കിഷ്ടമായി. ശിവഗിരിയിൽ സവർണരുടെ സത്യാഗ്രഹജാഥയ്ക്കു ഗുരു സ്വീകരണം ഏർപ്പാടു ചെയ്തു.
	
	ആര്യസമാജ പ്രവർത്ത‌കനായ സ്വാമി ശ്രദ്ധാനന്ദജി ശിവഗിരിയിൽ വന്ന് ഗുരുവിനെ അക്കൊല്ലം സന്ദർശിക്കുകയുണ്ടായി.
	
	1101 കന്നിയിൽ ഗുരു ബോധാനന്ദസ്വാമികളെ പിൻഗാമിയായി അഭിഷേകം ചെയ്തു. തുലാമാസം ബ്രഹ്മവിദ്യാലയത്തിന്റെ ശിലാസ്ഥാപനം നടത്തി. മേടത്തിൽ ഗുരു വിൽപ്പത്രം എഴുതി.
	
	അക്കൊല്ലം ദിവാൻ വാട്സ് ശിവഗിരി സന്ദർശിച്ചു. മുൻകൂട്ടി നിശ്ചയിച്ച സമയത്ത് ദിവാന് എത്തിച്ചേരാൻ കഴിഞ്ഞില്ല, റോഡിന്റെ ദുഃസ്ഥിതിയും മഴയുമാണ് കാരണമായി പറഞ്ഞത്. ``റോഡുകൾ ഇനി നന്നായിക്കൊള്ളുമല്ലോ" എന്നു ഗുരു പുഞ്ചിരിയോടെ പ്രതിവചിച്ചു.
	
	തലശ്ശേരി ജഗന്നാഥ ക്ഷേത്രത്തിനു മുന്നിൽ ഗുരുവിന്റെ ലോഹ പ്രതിമ സ്ഥാപിച്ചു. മുർക്കോത്തു കുമാരൻ അക്കാര്യത്തിൽ മുന്നിട്ടുനിന്നു പ്രവർത്തിച്ചു. ഒരു ഇറ്റാലിയൻ ശില്പി തീർത്ത ആ പ്രതിമ കണ്ടപ്പോൾ ഗുരു പറഞ്ഞു. ``ഇതു ജീവിച്ചുകൊള്ളും. ഭക്ഷണം ആവശ്യമില്ലല്ലൊ." അക്കൊല്ലം കുറ്റാലത്ത് സ്വാമികൾ കുറെ ദിവസം വിശ്രമിച്ചു.
	
	1102-ൽ സ്വാമി സത്യവ്രതൻ സമാധിയായി. സ്വാമികളെ ഈ സംഭവം സ്വല്പം ചലിപ്പിച്ചു. ``നമുക്കു ക്ഷീണം തോന്നുന്നു, സത്യവ്രതൻ സത്യവ്രതനായിത്തന്നെ ജീവിച്ചു, മരിച്ചു." എന്നു സ്വാമികൾ പറഞ്ഞു.
	
	കന്നിമാസം സ്വാമികൾ രണ്ടാം വട്ടം സിലോൺ സന്ദർശിച്ചു. മകരത്തിൽ സംഘടനാ സന്ദേശം പുറപ്പെടുവിച്ചു. ``സംഘടിച്ചു ശക്തരാവുക" എന്ന ആഹ്വാനം അന്നാണുണ്ടായത്. പള്ളാത്തുരുത്ത് എസ്. എൻ. ഡി. പി. യോഗവാർഷികത്തിൽ സ്വാമികൾ സന്നിഹിതനായി. എസ്. എൻ. ഡി. പി. യുടെ സംഘടനാപരമായ കാര്യങ്ങൾ ആലോചിക്കാൻ ശിവഗിരിയിൽ യോഗം ചേർന്നു. ഇടവത്തിൽ കളവങ്കോടം പ്രതിഷ്ഠ നടത്തി. ഇത് കണ്ണാടിപ്രതിഷ്ഠ എന്ന പേരിലാണ് പരക്കെ അറിയപ്പെടുന്നത്. ഓങ്കാര പ്രതിഷ്ഠ അല്ലെങ്കിൽ പ്രണവപ്രതിഷ്ഠ എന്നാണു സ്വാമികളുടെ സങ്കല്പം. കണ്ണാടിയിൽ ഓം ശാന്തി എന്നു കാണത്തക്ക വിധം പുറകിലെ രസം ചുരണ്ടിക്കളയാൻ പറഞ്ഞതിൽ നിന്നിതു വ്യക്തമാക്കുന്നു. വൈക്കത്തെ ഉല്ലലയിലെ കണ്ണാടിപ്രതിഷ്ഠ ഓങ്കാരേശ്വരപ്രതിഷ്ഠ എന്ന പേരിൽ അറിയുന്നു എന്നതും ഓർക്കാവുന്നതാണ്.
	
	1103 മകരത്തിൽ ശ്രീനാരായണ ധർമ്മസംഘം രജിസ്റ്റർ ചെയ്തു. കോട്ടയം എസ്. എൻ. ഡി. പി. യോഗത്തിൽ സംബന്ധിച്ചു. ശാഖായോഗങ്ങൾക്ക് സർട്ടിഫിക്കറ്റുകൾ വിതരണം ചെയ്തു. ശിവഗിരി തീർത്ഥാടനത്തിന് അനുമതി നൽകി. തീർത്ഥടനംകൊണ്ടു സാധിക്കേണ്ടതായി ഗുരു എണ്ണിപ്പറഞ്ഞ കാര്യങ്ങൾ താഴെ ചേർക്കുന്നു,
	
	1 വിദ്യാഭ്യാസം, 2 ശുചിത്വം 3 ഈശ്വരഭക്തി, 4 സംഘടന, 5 കൃഷി, 6 കച്ചവടം, 7 കൈത്തൊഴിൽ, 8 സാങ്കേതികശാസ്ത്രപരിശീലനങ്ങൾ.
	
	അതേ ആണ്ടിൽ ഗുരുദേവന് ശാരീരികാസ്വാസ്ഥ്യം ആരംഭിച്ചു. മൂത്രാശയ സംബന്ധമായ തകരാറായിരുന്നു രോഗ കാരണം. രോഗം, കൃത്യമായി നിർണയിക്കപ്പെടാഞ്ഞതുകൊണ്ടോ എന്തോ, ചികിത്സകൾ ഫലിക്കാതെ അതുവർദ്ധിച്ചു വന്നു. അലോപ്പതിയും ആയുർവേദവും പരീക്ഷിച്ചു. ചില ഘട്ടങ്ങളിൽ ഗുരുദേവൻ സ്വയം ചികിൽസ നടത്തുക പോലും ഉണ്ടായി. അതതു ചികിൽസാരീതികളിൽ ഏറ്റവും വിദഗ്ധരായവരുടെ സേവനം ലഭ്യമായി. പാലക്കാട്ടും മദ്രാസിലും പോയി ചികിൽസ നടത്തി. രോഗം കുറഞ്ഞില്ല.
	
	ശിവഗിരിയിൽതന്നെ മടങ്ങിവന്നു. ഗുരുദേവന്റെ രോഗവാർത്ത അറിഞ്ഞതു മുതൽ ജനങ്ങൾ ഉത്ക്കണ്ഠാകുലരായി കഴിഞ്ഞുകൂടി. ധർമ്മം പത്രത്തിൽ രോഗവിവരം അപ്പോഴപ്പോൾ പ്രസിദ്ധപ്പെടുത്തിക്കൊണ്ടിരുന്നു. 1104 കന്നി അഞ്ചിനു പകൽ മൂന്നു മണിയോടടുപ്പിച്ചു ആ ആദിമഹസ്സു ജഡമായ ശരീരം വെടിഞ്ഞു. ലക്ഷോപലക്ഷം ജനങ്ങളെ ഇരുളിൽ നിന്നു വെളിച്ചത്തിലേയ്ക്കു നയിച്ച ആ ആദിമഹസ്സ് അന്നെന്നപോലെ ഇന്നും വെളിച്ചം കാണിച്ചുകൊണ്ടിരിക്കുന്നു. അതെന്നും ആ കൃത്യം തുടർന്നുകൊണ്ടേയിരിക്കും.
	
	\newpage
	\thispagestyle{empty}
	\mbox{}
	
	\newpage
	\thispagestyle{empty}
	\mbox{}
	
	\newpage
	\thispagestyle{empty}
	
	\vspace*{\fill} % പേജിന്റെ മുകളിൽ നിന്ന് താഴേക്ക് തള്ളാൻ ഇത് സഹായിക്കും
	
	\begin{flushright}
		\large
		{\LARGE ശ്രീനാരായണ ഗുരു} \\
		ഗുരു ജീവിച്ചിരുന്നപ്പോൾ എഴുതപ്പെട്ട അദ്ദേഹത്തിന്റെ ഏക ജീവചരിത്രം \\[2\baselineskip]
		\textbf{ഒരു \href{URL}{വിക്കിഗ്രന്ഥശാല} പ്രസിദ്ധീകരണം}
	\end{flushright}
	
	\vspace*{2cm} % പേജിന്റെ ഏറ്റവും അടിയിൽ നിന്ന് അല്പം മുകളിലേക്ക് കയറ്റി നിർത്താൻ
	
\end{document}